% gpl-lgpl.tex                                                  -*- LaTeX -*-
%      Tutorial Text for the Detailed Study and Analysis of GPL and LGPL course
%
% Copyright (C) 2003, 2004, 2005 Free Software Foundation, Inc.

% License: CC-By-SA-4.0

% The copyright holders hereby grant the freedom to copy, modify, convey,
% Adapt, and/or redistribute this work under the terms of the Creative
% Commons Attribution Share Alike 4.0 International License.

% This text is distributed in the hope that it will be useful, but
% WITHOUT ANY WARRANTY; without even the implied warranty of
% MERCHANTABILITY or FITNESS FOR A PARTICULAR PURPOSE.

% You should have received a copy of the license with this document in
% a file called 'CC-By-SA-4.0.txt'.  If not, please visit
% https://creativecommons.org/licenses/by-sa/4.0/legalcode to receive
% the license text.

\documentclass[11pt, letterpaper]{book}
% FILTER_PS:  \usepackage[dvips]{graphicx}

\newcommand\href[2]{#2}

%\newcommand\textcolor[2]{#2}

\input{/usr/share/texmf/tex/plain/dvips/colordvi.tex}

% FILTER_PDF: \usepackage[pdftex]{graphicx}
\usepackage[
pdftex=true,
latex2html=false,
pdftitle={The GNU General Public License for Businesspeople and Developers},
pdfauthor={Bradley M. Kuhn},
pdfsubject={GNU General Public License},
pdfkeywords={computer, science, free, software, freedom, licensing, licenses, GPL, GNU, general, public, license}
]{hyperref}

% I could not get this to work!
%\usepackage[pdftex,usenames]{color}
% So I used this:

\input{/usr/share/texmf/pdftex/plain/misc/pdfcolor.tex}

%\usepackage[ref]{backref}

% FILTER_HTML: 
%\usepackage{html}

%\usepackage[dvips]{graphicx}

%\usepackage[latex2html]{hyperref}
\newcommand\href[2]{\htmladdnormallink{#1}{#2}}
%\renewcommand\cite[1]{\hypercite{#1}}
%\renewcommand\ref[1]{\hyperref{#1}}

\usepackage[usenames]{color}

%\input{/usr/share/texmf/tex/plain/dvips/colordvi.tex}

% NOT FOUND % one-inch-margins.tex                                            -*- LaTeX -*-
%      code to create one inch margins in LaTeX
%
% COPYRIGHT (C) 1994 Bradley M. Kuhn
%
% Written   :   Bradley M. Kuhn         Loyola College
%   By


%% This code creates one inch margins for a LaTeX document

\oddsidemargin 0in
\evensidemargin 0in
\textwidth 6.5in

\topmargin 0in
\textheight 8.0in

\usepackage{enumerate}
\usepackage[dvips]{graphicx}

%\setlength\parskip{0.7em}
%\setlength\parindent{0pt}

\newcommand{\defn}[1]{\emph{#1}}

%\pagestyle{empty}

\begin{document}

\frontmatter

\begin{titlepage}

\begin{center}

%\vspace{.5in}
\vfill

%\includegraphics{fsf-logo.eps}

\vfill

{\Large
{\sc The GPL and Legal Aspects of Free Software Development } \\

\vfill

%\vspace{.7in}



% \vspace{.3in}

Columbia Law School, New York, NY, USA \\
\vspace{.1in}
Wednesday, 28 September 2005
}

%\vspace{.7in}
\vfill

{\large
Daniel Ravicher

Senior Counsel 

Free Software Foundation

President and Executive Director

Public Patent Foundation
}

\vspace{.3in}


{\large
David Turner

GPL Compliance Engineer

Free Software Foundation
}

\end{center}

\vfill

{\parindent 0in
Copyright \copyright{} 2003, 2004, 2005 \hspace{.2in} Free Software Foundation, Inc.

\vspace{.3in}

The copyright holders hereby grant the freedom to copy, modify, convey,
Adapt, and/or redistribute this work under the terms of the Creative Commons
Attribution Share Alike 4.0 International License.  A copy of that license is
available at \verb=https://creativecommons.org/licenses/by-sa/4.0/legalcode=.
}

\end{titlepage}

\pagestyle{plain}

\pagenumbering{roman}

\part{The GPL and Legal Aspects of Free Software Development}

\chapter*{Preface}

This one-day course gives an explanation of the most popular Free
Software copyright license, the GNU General Public License (GNU GPL),
and teaches lawyers, software developers, managers and business people
how to use the GPL (and GPL'd software) successfully in a new Free
Software business and in existing, successful enterprises.

Attendees should have a general familiarity with software development
processes. A vague understanding of how copyright law applies to software
is also helpful. The tutorial is of most interest to lawyers, software
developers and managers who run software businesses that modify and/or
redistribute software under terms of the GNU GPL (or who wish to do so in
the future), and those who wish to make use of existing GPL'd software in
their enterprise.

Upon completion of the tutorial, successful attendees can expect to have
learned the following:

\begin{itemize}

  \item The freedom-defending purpose of each term of the GNU GPL

  \item The redistribution options under the GPL

  \item The obligations when modifying GPL'd software

  \item How to build a plan for proper and successful compliance with the GPL

  \item The business advantages that the GPL provides

  \item The most common business models used in conjunction with the GPL

  \item How existing GPL'd software can be used in existing enterprises

  \item The basics of the LGPL and how it differs from the GPL

  \item How best to understand the complexities regarding derivative
        works of software
\end{itemize}

\bigskip

These course materials are merely a summary of the highlights of the
course presented. Please be aware that during the actual GPL course, class
discussion supplements this printed curriculum. Simply reading it is
not equivalent to attending the course.

\tableofcontents

\pagebreak

%%%%%%%%%%%%%%%%%%%%%%%%%%%%%%%%%%%%%%%%%%%%%%%%%%%%%%%%%%%%%%%%%%%%%%%%%%%%%%%
% END OF ABSTRACTS SECTION
%%%%%%%%%%%%%%%%%%%%%%%%%%%%%%%%%%%%%%%%%%%%%%%%%%%%%%%%%%%%%%%%%%%%%%%%%%%%%%%
% START OF DAY ONE COURSE
%%%%%%%%%%%%%%%%%%%%%%%%%%%%%%%%%%%%%%%%%%%%%%%%%%%%%%%%%%%%%%%%%%%%%%%%%%%%%%%
\mainmatter

\pagenumbering{arabic}

\chapter{What Is Free Software?}

Consideration of the GNU General Public License (herein, abbreviated as
\defn{GNU GPL} or just \defn{GPL}) must begin by first considering the broader
world of Free Software. The GPL was not created from a void, rather,
it was created to embody and defend a set of principles that were set
forth at the founding of the GNU project and the Free Software Foundation
(FSF)---the organization that upholds, defends and promotes the philosophy
of software freedom. A prerequisite for understanding the GPL and its
terms and conditions is a basic understanding of the principles behind it.
The GPL is unlike almost all other software licenses in that it is
designed to defend and uphold these principles.

\section{The Free Software Definition}
\label{Free Software Definition}

The Free Software Definition is set forth in full on FSF's Web site at
\verb0http://www.fsf.org/0 \verb0philosophy/free-sw.html0. This section
presents an abbreviated version that will focus on the parts that are most
pertinent to the terms of the GPL\@.

A particular program is Free Software if it grants a particular user of
that program, the following freedoms:

\begin{itemize}

\item The freedom to run the program for any purpose

\item The freedom to change and modify the program

\item The freedom to copy and share the program

\item The freedom to share improved versions of the program

\end{itemize}

The focus on ``a particular user'' is very pertinent here. It is not
uncommon for the same version of a specific program to grant these
freedoms to some subset of its user base, while others have none or only
some of these freedoms. Section~\ref{Proprietary Relicensing} talks in
detail about how this can happen even if a program is released under the
GPL\@.

Some people refer to software that gives these freedoms as ``Open
Source.''  Besides having a different political focus than those who
call it Free Software,\footnote{The political differences between the
  Free Software Movement and the Open Source Movement are documented
  on FSF's Web site at {\tt
    http://www.fsf.org/licensing/essays/free-software-for-freedom.html}.}
those who call the software ``Open Source'' are focused on a side
issue.  User access to the source code of a program is a prerequisite
to make use of the freedom to modify. However, the important issue is
what freedoms are granted in the license of that source code.
Microsoft's ``Shared Source'' program, for example, gives various
types of access to source code, but almost none of the freedoms
described in this section.

One key issue central to these freedoms is that there are no
restrictions on how these freedoms can be exercised. Specifically, users
and programmers can exercise these freedoms noncommercially or
commercially. Licenses that grant these freedoms for noncommercial
activities but prohibit them for commercial activities are considered
non-Free.

In general, software for which most or all of these freedoms are
restricted in any way is called ``non-Free Software.''  Typically, the
term ``proprietary software'' is used more or less interchangeably with
``non-Free Software.''  Personally, I tend to use the term ``non-Free
Software'' to refer to noncommercial software that restricts freedom
(such as ``shareware'') and ``proprietary software'' to refer to
commercial software that restricts freedom (such as nearly all of
Microsoft's and Oracle's offerings).

The remainder of this section considers each of the four freedoms in
detail.

\subsection{The Freedom to Run}

For a program to be Free Software, the freedom to run the program must
be completely unrestricted. This means any use for software the user
can come up with must be permitted. Perhaps, for example, the user
has discovered an innovative use for a particular program, one
that the programmer never could have predicted. Such a use must not
be restricted.

It was once rare that this freedom was restricted by even proprietary
software; today it is not so rare. Most End User Licensing Agreements
(EULAs) that cover most proprietary software restrict some types of
use.  For example, some versions of Microsoft's FrontPage software
prohibit use of the software to create Web sites that generate
negative publicity for Microsoft. Free Software has no such
restrictions; everyone is free to use Free Software for any purpose
whatsoever.

\subsection{The Freedom to Change and Modify}

Free Software programs allow users to change, modify and adapt the
software to suit their needs. Access to the source code and related build
scripts are an essential part of this freedom. Without the source code
and the ability to build the binary applications from that source, the
freedom cannot be properly exercised.

Programmers can take direct benefit from this freedom, and often do.
However, this freedom is also important to users who are not programmers.
Users must have the right to exercise this freedom indirectly in both
commercial and noncommercial settings. For example, users often seek
noncommercial help with the software on email lists and in users groups.
When they find such help, they must have the freedom to recruit
programmers who might altruistically assist them to modify their software.

The commercial exercise of this freedom is also essential for users. Each
user, or group of users, must have the right to hire anyone they wish in a
competitive free market to modify and change the software. This means
that companies have a right to hire anyone they wish to modify their Free
Software. Additionally, such companies may contract with other companies
to commission software modification.

\subsection{The Freedom to Copy and Share}

Users may share Free Software in a variety of ways. Free Software
advocates work to eliminate a fundamental ethical dilemma of the software
age: choosing between obeying a software license, and friendship (by
giving away a copy of a program to your friend who likes the software you are
using). Free Software licenses, therefore, must permit this sort of
altruistic sharing of software among friends.

The commercial environment must also have the benefits of this freedom.
Commercial sharing typically takes the form of selling copies of Free
Software. Free Software can be sold at any price to anyone. Those who
redistribute Free Software commercially have the freedom to selectively
distribute (you can pick your customers) and to set prices at any level
the redistributor sees fit.

It is true that many people get copies of Free Software very cheaply (and
sometimes without charge). The competitive free market of Free Software
tends to keep prices low and reasonable. However, if someone is willing
to pay a billion dollars for one copy of the GNU Compiler Collection, such
a sale is completely permitted.

Another common instance of commercial sharing is service-oriented
distribution. For example, a distribution vendor may provide immediate
security and upgrade distribution via a special network service. Such
distribution is completely permitted for Free Software.

(Section~\ref{Business Models} of this tutorial talks in detail about
various Free Software business models that take advantage of the freedom
to share commercially.)

\subsection{The Freedom to Share Improvements}

The freedom to modify and improve is somewhat empty without the freedom to
share those improvements. The Free Software community is built on the
pillar of altruistic sharing of improved Free Software. Inevitably, a
Free Software project sprouts a mailing list where improvements are shared
freely among members of the development community. Such noncommercial
sharing must be permitted for Free Software to thrive.

Commercial sharing of modified Free Software is equally important.
For commercial support to exist in a competitive free market, all
developers --- from single-person contractors to large software
companies --- must have the freedom to market their services as
improvers of Free Software. All forms of such service marketing must
be equally available to all.

For example, selling support services for Free Software is fully
permitted. Companies and individuals can offer themselves as ``the place
to call'' when software fails or does not function properly. For such a
service to be meaningful, the entity offering that service must have the
right to modify and improve the software for the customer to correct any
problems that are beyond mere user error.

Entities must also be permitted to make available modified versions of
Free Software. Most Free Software programs have a ``standard version''
that is made available from the primary developers of the software.
However, all who have the software have the ``freedom to fork'' --- that
is, make available nontrivial modified versions of the software on a
permanent or semi-permanent basis. Such freedom is central to vibrant
developer and user interaction.

Companies and individuals have the right to make true value-added versions
of Free Software. They may use freedom to share improvements to
distribute distinct versions of Free Software with different functionality
and features. Furthermore, this freedom can be exercised to serve a
disenfranchised subset of the user community. If the developers of the
standard version refuse to serve the needs of some of the software's
users, other entities have the right to create a long- or short-lived fork
to serve that sub-community.

\section{How Does Software Become Free?}

The last section set forth the freedoms and rights respected by Free
Software. It presupposed, however, that such software exists. This
section discusses how Free Software comes into existence. But first, it
addresses how software can be non-Free in the first place.

Software can be made proprietary only because it is governed by copyright
law.\footnote{This statement is a bit of an oversimplification. Patents
  and trade secrets can cover software and make it effectively non-Free,
  one can contract away their rights and freedoms regarding software, or
  source code can be practically obscured in binary-only distribution
  without reliance on any legal system. However, the primary control
  mechanism for software is copyright.} Copyright law, with respect to
software, governs copying, modifying, and redistributing that
software.\footnote{Copyright law in general also governs ``public
  performance'' of copyrighted works. There is no generally agreed
  definition for public performance of software and version 2 of the GPL
  does not govern public performance.} By law, the copyright holder (a.k.a.
the author) of the work controls how others may copy, modify and/or
distribute the work. For proprietary software, these controls are used to
prohibit these activities. In addition, proprietary software distributors
further impede modification in a practical sense by distributing only
binary code and keeping the source code of the software secret.

Copyright law is a construction. In the USA, the Constitution permits,
but does not require, the creation of copyright law as federal
legislation. Software, since it is an idea fixed in a tangible medium, is
thus covered by the statues, and is copyrighted by default.

However, this legal construction is not necessarily natural. Software, in
its natural state without copyright, is Free Software. In an imaginary
world with no copyright, the rules would be different. In this
world, when you received a copy of a program's source code, there would be
no default legal system to restrict you from sharing it with others,
making modifications, or redistributing those modified
versions.\footnote{There could still exist legal systems, like our modern
  patent system, which could restrict the software in other ways.}

Software in the real world is copyrighted by default and is
automatically covered by that legal system. However, it is possible
to move software out of the domain of the copyright system. A
copyright holder is always permitted to \defn{disclaim} their
copyright. If copyright is disclaimed, the software is not governed
by copyright law. Software not governed by copyright is in the
``public domain.''

\subsection{Public Domain Software}
% FIXME: this section needs more improvements to make it clear that public
% domain dedication is difficult, if not impossible.
% Karen suggests that talking about USA government software being public
% domain might make sense here.
Theoretically, an author can create public domain software by disclaiming all
copyright interest on the work. In the USA and other countries that have
signed the Berne convention on copyright, software is copyrighted
automatically by the author when she ``fixes the software into a tangible
medium.''  In the software world, this usually means typing the source code
of the software into a file.

Imagine if an author can truly disclaim that default control given to her by the
copyright laws. Once this is done, the software is in the public domain
--- it is no longer covered by copyright. Since it is copyright law that
allows for various controls on software (i.e., prohibition of copying,
modification, and redistribution), removing the software from the
copyright system and placing it into the public domain does yield Free
Software.

Carefully note that software in the public domain is \emph{not} licensed
in any way. It is nonsensical to say software is ``licensed for the
public domain,'' or any phrase that implies the copyright holder gave
expressed permission to take actions governed by copyright law.

By contrast, what the copyright holder has done is renounce her copyright
controls on the work. The law gave her controls over the work, and she
has chosen to waive those controls. Software in the public domain is
absent copyright and absent a license. The software freedoms discussed in
Section~\ref{Free Software Definition} are all granted because there is no
legal system in play to take them away.

\subsection{Why Copyright Free Software?}

If simply disclaiming copyright on software yields Free Software, then it
stands to reason that putting software into the public domain is the
easiest and most straightforward way to produce Free Software. Indeed,
some major Free Software projects have chosen this method for making their
software Free. However, most of the Free Software in existence \emph{is}
copyrighted. In most cases (particularly in those of FSF and the GNU
Project), this was done due to very careful planning.

Software released into the public domain does grant freedom to those users
who receive the standard versions on which the original author disclaimed
copyright. However, since the work is not copyrighted, any nontrivial
modification made to the work is fully copyrightable.

Free Software released into the public domain initially is Free, and
perhaps some who modify the software choose to place their work into the
public domain as well. However, over time, some entities will choose to
proprietarize their modified versions. The public domain body of software
feeds the proprietary software. The public commons disappears, because
fewer and fewer entities have an incentive to contribute back to the
commons. They know that any of their competitors can proprietarize their
enhancements. Over time, almost no interesting work is left in the public
domain, because nearly all new work is done by proprietarization.

A legal mechanism is needed to redress this problem. FSF was in fact
originally created primarily as a legal entity to defend software freedom,
and that work of defending software freedom is a substantial part of
its work today. Specifically because of this ``embrace, proprietarize and
extend'' cycle, FSF made a conscious choice to copyright its Free Software,
and then license it under ``copyleft'' terms. Many, including the
developers of the kernel named Linux, have chosen to follow this paradigm.

Copyleft is a legal strategy to defend, uphold and propagate software
freedom. The basic technique of copyleft is as follows: copyright the
software, license it under terms that give all the software freedoms, but
use the copyright law controls to ensure that all who receive a copy of
the software have equal rights and freedom. In essence, copyleft grants
freedom, but forbids others to forbid that freedom to anyone else along
the distribution and modification chains.

Copyleft is a general concept. Much like ideas for what a computer might
do must be \emph{implemented} by a program that actually does the job, so
too must copyleft be implemented in some concrete legal structure.
``Share and share alike'' is a phrase that is used often enough to explain the
concept behind copyleft, but to actually make it work in the real world, a
true implementation in legal text must exist. The GPL is the primary
implementation of copyleft in copyright licensing language.

\section{A Community of Equality}

The GPL uses copyright law to defend freedom and equally ensure users'
rights. This ultimately creates an community of equality for both
business and noncommercial users.

\subsection{The Noncommercial Community}

A GPL'd code base becomes a center of a vibrant development and user
community. Traditionally, volunteers, operating noncommercially out of
keen interest or ``scratch an itch'' motivations, produce initial versions
of a GPL'd system. Because of the efficient distribution channels of the
Internet, any useful GPL'd system is adopted quickly by noncommercial
users.

Fundamentally, the early release and quick distribution of the software
gives birth to a thriving noncommercial community. Users and developers
begin sharing bug reports and bug fixes across a shared intellectual
commons. Users can trust the developers, because they know that if the
developers fail to address their needs or abandon the project, the GPL
ensures that someone else has the right to pick up development.
Developers know that the users cannot redistribute their software without
passing along the rights granted by GPL, so they are assured that every
one of their users is treated equally.

Because of the symmetry and fairness inherent in GPL'd distribution,
nearly every GPL'd package in existence has a vibrant noncommercial user
and developer base.

\subsection{The Commercial Community}

By the same token, nearly all established GPL'd software systems have a
vibrant commercial community. Nearly every GPL'd system that has gained
wide adoption from noncommercial users and developers eventually begins
to fuel a commercial system around that software.

For example, consider the Samba file server system that allows Unix-like
systems (including GNU/Linux) to serve files to Microsoft Windows systems.
Two graduate students originally developed Samba in their spare time and
it was deployed noncommercially in academic environments. However, very
soon for-profit companies discovered that the software could work for them
as well, and their system administrators began to use it in place of
Microsoft Windows NT file-servers. This served to lower the cost of
running such servers by orders of magnitude. There was suddenly room in
Windows file-server budgets to hire contractors to improve Samba. Some of
the first people hired to do such work were those same two graduate
students who originally developed the software.

The noncommercial users, however, were not concerned when these two
fellows began collecting paychecks off of their GPL'd work. They knew
that because of the nature of the GPL that improvements that were
distributed in the commercial environment could easily be folded back into
the standard version. Companies are not permitted to proprietarize
Samba, so the noncommercial users, and even other commercial users are
safe in the knowledge that the software freedom ensured by GPL will remain
protected.

Commercial developers also work in concert with noncommercial
developers. Those two now-long-since graduated students continue to
contribute to Samba altruistically, but also get paid work doing it.
Priorities change when a client is in the mix, but all the code is
contributed back to the standard version. Meanwhile, many other
individuals have gotten involved noncommercially as developers,
because they want to ``cut their teeth on Free Software,'' or because
the problems interest them. When they get good at it, perhaps they
will move on to another project, or perhaps they will become
commercial developers of the software themselves.

No party is a threat to another in the GPL software scenario because
everyone is on equal ground. The GPL protects rights of the commercial
and noncommercial contributors and users equally. The GPL creates trust,
because it is a level playing field for all.

\subsection{Law Analogy}

In his introduction to Stallman's \emph{Free Software, Free Society},
Lawrence Lessig draws an interesting analogy between the law and Free
Software. He argues that the laws of a free society must be protected
much like the GPL protects software. So that I might do true justice to
Lessig's argument, I quote it verbatim:

\begin{quotation}

A ``free society'' is regulated by law. But there are limits that any free
society places on this regulation through law: No society that kept its
laws secret could ever be called free. No government that hid its
regulations from the regulated could ever stand in our tradition. Law
controls. But it does so justly only when visibly. And law is visible
only when its terms are knowable and controllable by those it regulates,
or by the agents of those it regulates (lawyers, legislatures).

This condition on law extends beyond the work of a legislature. Think
about the practice of law in American courts. Lawyers are hired by their
clients to advance their clients' interests. Sometimes that interest is
advanced through litigation. In the course of this litigation, lawyers
write briefs. These briefs in turn affect opinions written by judges.
These opinions decide who wins a particular case, or whether a certain law
can stand consistently with a constitution.

All the material in this process is free in the sense that Stallman means.
Legal briefs are open and free for others to use. The arguments are
transparent (which is different from saying they are good), and the
reasoning can be taken without the permission of the original lawyers.
The opinions they produce can be quoted in later briefs. They can be
copied and integrated into another brief or opinion. The ``source code''
for American law is by design, and by principle, open and free for anyone
to take. And take lawyers do---for it is a measure of a great brief that
it achieves its creativity through the reuse of what happened before. The
source is free; creativity and an economy is built upon it.

This economy of free code (and here I mean free legal code) doesn't starve
lawyers. Law firms have enough incentive to produce great briefs even
though the stuff they build can be taken and copied by anyone else. The
lawyer is a craftsman; his or her product is public. Yet the crafting is
not charity. Lawyers get paid; the public doesn't demand such work
without price. Instead this economy flourishes, with later work added to
the earlier.

We could imagine a legal practice that was different---briefs and
arguments that were kept secret; rulings that announced a result but not
the reasoning. Laws that were kept by the police but published to no one
else. Regulation that operated without explaining its rule.

We could imagine this society, but we could not imagine calling it
``free.''  Whether or not the incentives in such a society would be better
or more efficiently allocated, such a society could not be known as free.
The ideals of freedom, of life within a free society, demand more than
efficient application. Instead, openness and transparency are the
constraints within which a legal system gets built, not options to be
added if convenient to the leaders. Life governed by software code should
be no less.

Code writing is not litigation. It is better, richer, more
productive. But the law is an obvious instance of how creativity and
incentives do not depend upon perfect control over the products
created. Like jazz, or novels, or architecture, the law gets built
upon the work that went before. This adding and changing is what
creativity always is. And a free society is one that assures that its
most important resources remain free in just this sense.\footnote{This
quotation is Copyright \copyright{} 2002, Lawrence Lessig. It is
licensed under the terms of
\texttt{http://creativecommons.org/licenses/by/1.0/}{the ``Attribution
License'' version 1.0} or any later version as published by Creative
Commons.}
\end{quotation}

In essence, lawyers are paid to service the shared commons of legal
infrastructure. Few citizens defend themselves in court or write their
own briefs (even though they are legally permitted to do so) because
everyone would prefer to have an expert do that job.

The Free Software economy is a market ripe for experts. It
functions similarly to other well established professional fields like the
law. The GPL, in turn, serves as the legal scaffolding that permits the
creation of this vibrant commercial and noncommercial Free Software
economy.

%%%%%%%%%%%%%%%%%%%%%%%%%%%%%%%%%%%%%%%%%%%%%%%%%%%%%%%%%%%%%%%%%%%%%%%%%%%%%%%
\chapter{Running Software and Verbatim Copying}
\label{run-and-verbatim}


This chapter begins the deep discussion of the details of the terms of
GPL\@. In this chapter, we consider the first two sections: GPL \S\S
0--2. These are the straightforward sections of the GPL that define the
simplest rights that the user receives.

\section{GPL \S 0: Freedom to Run}
\label{GPLs0}

\S 0, the opening section of GPL, sets forth that the work is governed by
copyright law. It specifically points out that it is the ``copyright
holder'' who decides if a work is licensed under its terms and explains
how the copyright holder might indicate this fact.

A bit more subtly, \S 0 makes an inference that copyright law is the only
system under which it is governed. Specifically, it states:
\begin{quote}
Activities other than copying, distribution and modification are not
covered by this License; they are outside its scope.
\end{quote}
In essence, the license governs \emph{only} those activities, and all other
activities are unrestricted, provided that no other agreements trump GPL
(which they cannot; see Sections~\ref{GPLs6} and~\ref{GPLs7}). This is
very important, because the Free Software community heavily supports
users' rights to ``fair use'' and ``unregulated use'' of copyrighted
material. GPL asserts through this clause that it supports users' rights
to fair and unregulated uses.

Fair use of copyrighted material is an established legal doctrine that
permits certain activities. Discussion of the various types of fair
use activity are beyond the scope of this tutorial. However, one
important example of fair use is the right to quote a very few lines
(less than seven or so) and reuse them as you would with or without
licensing restrictions.

Fair use is a doctrine established by the courts or by statute. By
contrast, unregulated uses are those that are not covered by the statue
nor determined by a court to be covered, but are common and enjoyed by
many users. An example of unregulated use is reading a printout of the
program's source code like an instruction book for the purpose of learning
how to be a better programmer.

\medskip

Thus, the GPL protects users fair and unregulated use rights precisely by
not attempting to cover them. Furthermore, the GPL ensures the freedom
to run specifically by stating the following:
\begin{quote}
''The act of running the Program is not restricted.''
\end{quote}
Thus, users are explicitly given the freedom to run by \S 0.

\medskip

The bulk of \S 0 not yet discussed gives definitions for other terms used
throughout. The only one worth discussing in detail is ``work based on
the Program.''  The reason this definition is particularly interesting is
not for the definition itself, which is rather straightforward, but
because it clears up a common misconception about the GPL\@.

The GPL is often mistakenly criticized because it fails to give a
definition of ``derivative work.''  In fact, it would be incorrect and
problematic if the GPL attempted to define this. A copyright license, in
fact, has no control over what may or may not be a derivative work. This
matter is left up to copyright law, not the licenses that utilize it.

It is certainly true that copyright law as a whole does not propose clear
and straightforward guidelines for what is and is not a derivative
software work under copyright law. However, no copyright license --- not
even the GNU GPL --- can be blamed for this. Legislators and court
opinions must give us guidance to decide the border cases.

\section{GPL \S 1: Verbatim Copying}
\label{GPLs1}

GPL \S 1 covers the matter of redistributing the source code of a program
exactly as it was received. This section is quite straightforward.
However, there are a few details worth noting here.

The phrase ``in any medium'' is important. This, for example, gives the
freedom to publish a book that is the printed copy of the program's source
code. It also allows for changes in the medium of distribution. Some
vendors may ship Free Software on a CD, but others may place it right on
the hard drive of a pre-installed computer. Any such redistribution media
is allowed.

Preservation of copyright notice and license notifications are mentioned
specifically in \S 1. These are in some ways the most important part of
the redistribution, which is why they are mentioned by name. The GPL
always strives to make it abundantly clear to anyone who receives the
software what its license is. The goal is to make sure users know their
rights and freedoms under GPL, and to leave no reason that someone would be
surprised the software she got was licensed under GPL\@. Thus
throughout the GPL, there are specific references to the importance of
notifying others down the distribution chain that they have rights under
GPL.

Also mentioned by name is the warranty disclaimer. Most people today do
not believe that software comes with any warranty. Notwithstanding the
proposed state-level UCITA bills (which have never obtained widespread
adoption), there are few or no implied warranties with software.
However, just to be on the safe side, GPL clearly disclaims them, and the
GPL requires redistributors to keep the disclaimer very visible. (See
Sections~\ref{GPLs11} and~\ref{GPLs12} of this tutorial for more on GPL's
warranty disclaimers.)

Note finally that \S 1 begins to set forth the important defense of
commercial freedom. \S 1 clearly states that in the case of verbatim
copies, one may make money. Redistributors are fully permitted to charge
for the redistribution of copies of Free Software. In addition, they may
provide the warranty protection that the GPL disclaims as an additional
service for a fee. (See Section~\ref{Business Models} for more discussion
on making a profit from Free Software redistribution.)

%%%%%%%%%%%%%%%%%%%%%%%%%%%%%%%%%%%%%%%%%%%%%%%%%%%%%%%%%%%%%%%%%%%%%%%%%%%%%%%

\chapter{Derivative Works: Statute and Case Law}

We digress for this chapter from our discussion of GPL's exact text to
consider the matter of derivative works --- a concept that we must
understand fully before considering \S\S 2--3 of GPL\@. GPL, and Free
Software licensing in general, relies critically on the concept of
``derivative work'' since software that is ``independent,'' (i.e., not
``derivative'') of Free Software need not abide by the terms of the
applicable Free Software license. As much is required by \S 106 of the
Copyright Act, 17 U.S.C. \S 106 (2002), and admitted by Free Software
licenses, such as the GPL, which (as we have seen) states in \S 0 that ``a
`work based on the Program' means either the Program or any derivative
work under copyright law.'' It is being a derivative work of Free Software
that triggers the necessity to comply with the terms of the Free Software
license under which the original work is distributed. Therefore, one is
left to ask, just what is a ``derivative work''? The answer to that
question differs depending on which court is being asked.

The analysis in this chapter sets forth the differing definitions of
derivative work by the circuit courts. The broadest and most
established definition of derivative work for software is the
abstraction, filtration, and comparison test (``the AFC test'') as
created and developed by the Second Circuit. Some circuits, including
the Ninth Circuit and the First Circuit, have either adopted narrower
versions of the AFC test or have expressly rejected the AFC test in
favor of a narrower standard. Further, several other circuits have yet
to adopt any definition of derivative work for software.

As an introductory matter, it is important to note that literal copying of
a significant portion of source code is not always sufficient to establish
that a second work is a derivative work of an original
program. Conversely, a second work can be a derivative work of an original
program even though absolutely no copying of the literal source code of
the original program has been made. This is the case because copyright
protection does not always extend to all portions of a program's code,
while, at the same time, it can extend beyond the literal code of a
program to its non-literal aspects, such as its architecture, structure,
sequence, organization, operational modules, and computer-user interface.

\section{The Copyright Act}

The copyright act is of little, if any, help in determining the definition
of a derivative work of software. However, the applicable provisions do
provide some, albeit quite cursory, guidance. Section 101 of the Copyright
Act sets forth the following definitions:

\begin{quotation}
A ``computer program'' is a set of statements or instructions to be used
directly or indirectly in a computer in order to bring about a certain
result.

A ``derivative work'' is a work based upon one or more preexisting works,
such as a translation, musical arrangement, dramatization,
fictionalization, motion picture version, sound recording, art
reproduction, abridgment, condensation, or any other form in which a work
may be recast, transformed, or adapted. A work consisting of editorial
revisions, annotations, elaborations, or other modifications which, as a
whole, represent an original work of authorship, is a ``derivative work.''
\end{quotation}

These are the only provisions in the Copyright Act relevant to the
determination of what constitutes a derivative work of a computer
program. Another provision of the Copyright Act that is also relevant to
the definition of derivative work is \S 102(b), which reads as follows:

\begin{quotation}
In no case does copyright protection for an original work of authorship
extend to any idea, procedure, process, system, method of operation,
concept, principle, or discovery, regardless of the form in which it is
described, explained, illustrated, or embodied in such work.
\end{quotation}

Therefore, before a court can ask whether one program is a derivative work
of another program, it must be careful not to extend copyright protection
to any ideas, procedures, processes, systems, methods of operation,
concepts, principles, or discoveries contained in the original program. It
is the implementation of this requirement to ``strip out'' unprotectable
elements that serves as the most frequent issue over which courts
disagree.

\section{Abstraction, Filtration, Comparison Test}

As mentioned above, the AFC test for determining whether a computer
program is a derivative work of an earlier program was created by the
Second Circuit and has since been adopted in the Fifth, Tenth, and
Eleventh Circuits. Computer Associates Intl., Inc. v. Altai, Inc., 982
F.2d 693 (2nd Cir. 1992); Engineering Dynamics, Inc. v. Structural
Software, Inc., 26 F.3d 1335 (5th Cir. 1994); Kepner-Tregoe,
Inc. v. Leadership Software, Inc., 12 F.3d 527 (5th Cir. 1994); Gates
Rubber Co. v. Bando Chem. Indust., Ltd., 9 F.3d 823 (10th Cir. 1993);
Mitel, Inc. v. Iqtel, Inc., 124 F.3d 1366 (10th Cir. 1997); 5 Bateman
v. Mnemonics, Inc., 79 F.3d 1532 (11th Cir. 1996); and, Mitek Holdings,
Inc. v. Arce Engineering Co., Inc., 89 F.3d 1548 (11th Cir. 1996).

Under the AFC test, a court first abstracts from the original program its
constituent structural parts. Then, the court filters from those
structural parts all unprotectable portions, including incorporated ideas,
expression that is necessarily incidental to those ideas, and elements
that are taken from the public domain. Finally, the court compares any and
all remaining kernels of creative expression to the structure of the
second program to determine whether the software programs at issue are
substantially similar so as to warrant a finding that one is the
derivative work of the other.

Often, the courts that apply the AFC test will perform a quick initial
comparison between the entirety of the two programs at issue in order to
help determine whether one is a derivative work of the other. Such a
holistic comparison, although not a substitute for the full application of
the AFC test, sometimes reveals a pattern of copying that is not otherwise
obvious from the application of the AFC test when, as discussed below,
only certain components of the original program are compared to the second
program. If such a pattern is revealed by the quick initial comparison,
the court is more likely to conclude that the second work is indeed a
derivative of the original.

\subsection{Abstraction}

The first step courts perform under the AFC test is separation of the
work's ideas from its expression. In a process akin to reverse
engineering, the courts dissect the original program to isolate each level
of abstraction contained within it. Courts have stated that the
abstractions step is particularly well suited for computer programs
because it breaks down software in a way that mirrors the way it is
typically created. However, the courts have also indicated that this step
of the AFC test requires substantial guidance from experts, because it is
extremely fact and situation specific.

By way of example, one set of abstraction levels is, in descending order
of generality, as follows: the main purpose, system architecture, abstract
data types, algorithms and data structures, source code, and object
code. As this set of abstraction levels shows, during the abstraction step
of the AFC test, the literal elements of the computer program, namely the
source and object code, are defined as particular levels of
abstraction. Further, the source and object code elements of a program are
not the only elements capable of forming the basis for a finding that a
second work is a derivative of the program. In some cases, in order to
avoid a lengthy factual inquiry by the court, the owner of the copyright in
the original work will submit its own list of what it believes to be the
protected elements of the original program. In those situations, the court
will forgo performing its own abstraction, and proceed to the second step of
the AFC test.

\subsection{Filtration}

The most difficult and controversial part of the AFC test is the second
step, which entails the filtration of protectable expression contained in
the original program from any unprotectable elements nestled therein. In
determining which elements of a program are unprotectable, courts employ a
myriad of rules and procedures to sift from a program all the portions
that are not eligible for copyright protection.

First, as set forth in \S 102(b) of the Copyright Act, any and all ideas
embodied in the program are to be denied copyright protection. However,
implementing this rule is not as easy as it first appears. The courts
readily recognize the intrinsic difficulty in distinguishing between ideas
and expression and that, given the varying nature of computer programs,
doing so will be done on an ad hoc basis. The first step of the AFC test,
the abstraction, exists precisely to assist in this endeavor by helping
the court separate out all the individual elements of the program so that
they can be independently analyzed for their expressive nature.

A second rule applied by the courts in performing the filtration step of
the AFC test is the doctrine of merger, which denies copyright protection
to expression necessarily incidental to the idea being expressed. The
reasoning behind this doctrine is that when there is only one way to
express an idea, the idea and the expression merge, meaning that the
expression cannot receive copyright protection due to the bar on copyright
protection extending to ideas. In applying this doctrine, a court will ask
whether the program's use of particular code or structure is necessary for
the efficient implementation of a certain function or process. If so, then
that particular code or structure is not protected by copyright and, as a
result, it is filtered away from the remaining protectable expression.

A third rule applied by the courts in performing the filtration step of
the AFC test is the doctrine of scenes a faire, which denies copyright
protection to elements of a computer program that are dictated by external
factors. Such external factors can include:

\begin{itemize}

  \item The mechanical
specifications of the computer on which a particular program is intended
to operate

  \item Compatibility requirements of other programs with which a
program is designed to operate in conjunction

  \item Computer manufacturers'
design standards

  \item Demands of the industry being serviced, and

widely accepted programming practices within the computer industry

\end{itemize}

Any code or structure of a program that was shaped predominantly in
response to these factors is filtered out and not protected by
copyright. Lastly, elements of a computer program are also to be filtered
out if they were taken from the public domain or fail to have sufficient
originality to merit copyright protection.

Portions of the source or object code of a computer program are rarely
filtered out as unprotectable elements. However, some distinct parts of
source and object code have been found unprotectable. For example,
constant s, the invariable integers comprising part of formulas used to
perform calculations in a program, are unprotectable. Further, although
common errors found in two programs can provide strong evidence of
copying, they are not afforded any copyright protection over and above the
protection given to the expression containing them.

\subsection{Comparison}

The third and final step of the AFC test entails a comparison of the
original program's remaining protectable expression to a second
program. The issue will be whether any of the protected expression is
copied in the second program and, if so, what relative importance the
copied portion has with respect to the original program overall. The
ultimate inquiry is whether there is ``substantial'' similarity between
the protected elements of the original program and the potentially
derivative work. The courts admit that this process is primarily
qualitative rather than quantitative and is performed on a case-by-case
basis. In essence, the comparison is an ad hoc determination of whether
the protectable elements of the original program that are contained in the
second work are significant or important parts of the original program. If
so, then the second work is a derivative work of the first. If, however,
the amount of protectable elements copied in the second work are so small
as to be de minimis, then the second work is not a derivative work of the
original.

\section{Analytic Dissection Test}

The Ninth Circuit has adopted the analytic dissection test to determine
whether one program is a derivative work of another. Apple Computer,
Inc. v. Microsoft Corp., 35 F.3d 1435 (9th Cir. 1994). The analytic
dissection test first considers whether there are substantial similarities
in both the ideas and expressions of the two works at issue. Once the
similar features are identified, analytic dissection is used to determine
whether any of those similar features are protected by copyright. This
step is the same as the filtration step in the AFC test. After identifying
the copyrightable similar features of the works, the court then decides
whether those features are entitled to ``broad'' or ``thin''
protection. ``Thin'' protection is given to non-copyrightable facts or
ideas that are combined in a way that affords copyright protection only
from their alignment and presentation, while ``broad'' protection is given
to copyrightable expression itself. Depending on the degree of protection
afforded, the court then sets the appropriate standard for a subjective
comparison of the works to determine whether, as a whole, they are
sufficiently similar to support a finding that one is a derivative work of
the other. ``Thin'' protection requires the second work be virtually
identical in order to be held a derivative work of an original, while
``broad'' protection requires only a ``substantial similarity.''

\section{No Protection for ``Methods of Operation''}

The First Circuit expressly rejected the AFC test and, instead, takes a
much narrower view of the meaning of derivative work for software. The
First Circuit holds that ``method of operation,'' as used in \S 102(b) of
the Copyright Act, refers to the means by which users operate
computers. Lotus Development Corp. v. Borland Int’l., Inc., 49 F.3d 807
(1st Cir. 1995). More specifically, the court held that a menu command
hierarchy for a computer program was uncopyrightable because it did not
merely explain and present the program’s functional capabilities to the
user, but also served as a method by which the program was operated and
controlled. As a result, under the First Circuit’s test, literal copying
of a menu command hierarchy, or any other ``method of operation,'' cannot
form the basis for a determination that one work is a derivative of
another. It is also reasonable to expect that the First Circuit will read
the unprotectable elements set forth in \S 102(b) broadly, and, as such,
promulgate a definition of derivative work that is much narrower than that
which exists under the AFC test.

\section{No Test Yet Adopted}

Several circuits, most notably the Fourth and Seventh, have yet to
declare their definition of derivative work and whether or not the
AFC, Analytic Dissection, or some other test best fits their
interpretation of copyright law. Therefore, uncertainty exists with
respect to determining the extent to which a software program is a
derivative work of another in those circuits. However, one may presume
that they would give deference to the AFC test since it is by far the
majority rule amongst those circuits that have a standard for defining
a software derivative work.

\section{Cases Applying Software Derivative Work Analysis}

In the preeminent case regarding the definition of a derivative work for
software, Computer Associates v. Altai, the plaintiff alleged that its
program, Adapter, which was used to handle the differences in operating
system calls and services, was infringed by the defendant's competitive
program, Oscar. About 30\% of Oscar was literally the same code as
that in Adapter. After the suit began, the defendant rewrote those
portions of Oscar that contained Adapter code in order to produce a new
version of Oscar that was functionally competitive with Adapter, without
have any literal copies of its code. Feeling slighted still, the
plaintiff alleged that even the second version of Oscar, despite having no
literally copied code, also infringed its copyrights. In addressing that
question, the Second Circuit promulgated the AFC test.

In abstracting the various levels of the program, the court noted a
similarity between the two programs' parameter lists and macros. However,
following the filtration step of the AFC test, only a handful of the lists
and macros were protectable under copyright law because they were either
in the public domain or required by functional demands on the
program. With respect to the handful of parameter lists and macros that
did qualify for copyright protection, after performing the comparison step
of the AFC test, it was reasonable for the district court to conclude that
they did not warrant a finding of infringement given their relatively minor
contribution to the program as a whole. Likewise, the similarity between
the organizational charts of the two programs was not substantial enough
to support a finding of infringement because they were too simple and
obvious to contain any original expression.

Perhaps not surprisingly, there have been few cases involving a highly
detailed software derivative work analysis. Most often, cases involve
clearer basis for decision, including frequent bad faith on the part of
the defendant or overaggressiveness on the part of the plaintiff.
However, no cases involving Free Software licensing have ever gone to
court. As Free Software becomes an ever-increasingly important part of
the economy, it remains to be seen if battle lines will be
drawn over whether particular programs infringe the rights of Free
Software developers or whether the entire community, including industry,
adopts norms avoiding such risk.

%%%%%%%%%%%%%%%%%%%%%%%%%%%%%%%%%%%%%%%%%%%%%%%%%%%%%%%%%%%%%%%%%%%%%%%%%%%%%%%

\chapter{Modified Source and Binary Distribution}
\label{source-and-binary}

In this chapter, we discuss the two core sections that define the rights
and obligations for those who modify, improve, and/or redistribute GPL'd
software. These sections, \S\S 2--3, define the central core rights and
requirements of GPL\@.

\section{GPL \S 2: Share and Share Alike}

For many, this is where the ``magic'' happens that defends software
freedom along the distribution chain. \S 2 is the only place in the GPL
that governs the modification controls of copyright law. If someone
modifies a GPL'd program, she is bound in the making those changes by \S
2. The goal here is to ensure that the body of GPL'd software, as it
continues and develops, remains Free as in freedom.

To achieve that goal, \S 2 first sets forth that the rights of
redistribution of modified versions are the same as those for verbatim
copying, as presented in \S 1. Therefore, the details of charging,
keeping copyright notices intact, and other \S 1 provisions are in tact
here as well. However, there are three additional requirements.

The first (\S 2(a)) requires that modified files carry ``prominent
notices'' explaining what changes were made and the date of such
changes. The goal here is not to put forward some specific way of
marking changes nor controlling the process of how changes get made.
Primarily, \S 2(a) seeks to ensure that those receiving modified
versions know the history of changes to the software. For some users,
it is important to know that they are using the standard version of
program, because while there are many advantages to using a fork,
there are a few disadvantages. Users should be informed about the
historical context of the software version they use, so that they can
make proper support choices. Finally, \S 2(a) serves an academic
purpose --- ensuring that future developers can use a diachronic
approach to understand the software.

\medskip

The second requirement (\S 2(b)) contains the four short lines that embody
the legal details of ``share and share alike.''  These 46 words are
considered by some to be the most worthy of careful scrutiny because \S
2(b) can be a source of great confusion when not properly understood.

In considering \S 2(b), first note the qualifier: it only applies to
derivative works that ``you distribute or publish.''  Despite years of
education efforts by FSF on this matter, many still believe that modifiers
of GPL'd software are required by the license to publish or otherwise
share their changes. On the contrary, \S 2(b) {\bf does not apply if} the
changes are never distributed. Indeed, the freedom to make private,
personal, unshared changes to software for personal use only should be
protected and defended.\footnote{FSF does maintain that there is an {\bf
    ethical} obligation to redistribute changes that are generally useful,
  and often encourages companies and individuals to do so. However, there
  is a clear distinction between what one {\bf ought} to do and what one
  {\bf must} do.}

Next, we again encounter the same matter that appears in \S 0, in the
following text:
\begin{quote}
``...that in whole or part contains or is derived from the Program or any part thereof.''
\end{quote}
Again, the GPL relies here on what the copyright law says is a derivative
work. If, under copyright law, the modified version ``contains or is
derived from'' the GPL'd software, then the requirements of \S 2(b)
apply. The GPL invokes its control as a copyright license over the
modification of the work in combination with its control over distribution
of the work.

The final clause of \S 2(b) describes what the licensee must do if she is
distributing or publishing a work that is deemed a derivative work under
copyright law --- namely, the following:
\begin{quote}
[The work must] be licensed as a whole at no charge to all third parties
under the terms of this License.
\end{quote}
That is probably the most tightly-packed phrase in all of the GPL\@.
Consider each subpart carefully.

The work ``as a whole'' is what is to be licensed. This is an important
point that \S 2 spends an entire paragraph explaining; thus this phrase is
worthy of a lengthy discussion here. As a programmer modifies a software
program, she generates new copyrighted material --- fixing expressions of
ideas into the tangible medium of electronic file storage. That
programmer is indeed the copyright holder of those new changes. However,
those changes are part and parcel to the original work distributed to
the programmer under GPL\@. Thus, the license of the original work
affects the license of the new whole derivative work.

% {\cal I}
\newcommand{\gplusi}{$\mathcal{G\!\!+\!\!I}$}
\newcommand{\worki}{$\mathcal{I}$}
\newcommand{\workg}{$\mathcal{G}$}

\label{separate-and-independent}

It is certainly possible to take an existing independent work (called
\worki{}) and combine it with a GPL'd program (called \workg{}). The
license of \worki{}, when it is distributed as a separate and independent
work, remains the prerogative of the copyright holder of \worki{}.
However, when \worki{} is combined with \workg{}, it produces a new work
that is the combination of the two (called \gplusi{}). The copyright of
this combined work, \gplusi{}, is held by the original copyright
holder of each of the two works.

In this case, \S 2 lays out the terms by which \gplusi{} may be
distributed and copied. By default, under copyright law, the copyright
holder of \worki{} would not have been permitted to distribute \gplusi{};
copyright law forbids it without the expressed permission of the copyright
holder of \workg{}. (Imagine, for a moment, if \workg{} were a Microsoft
product --- would they give you permission to create and distribute
\gplusi{} without paying them a hefty sum?)  The license of \workg{}, the
GPL, sets forth ahead of time options for the copyright holder of \worki{}
who may want to create and distribute \gplusi{}. This pregranted
permission to create and distribute derivative works, provided the terms
of GPL are upheld, goes far above and beyond the permissions that one
would get with a typical work not covered by a copyleft license. Thus, to
say that this restriction is any way unreasonable is simply ludicrous.

\medskip

The next phrase of note in \S 2(b) is ``licensed...at no charge.''
This is a source of great confusion to many. Not a month goes by that
FSF does not receive an email that claims to point out ``a
contradiction in GPL'' because \S 2 says that redistributors cannot
charge for modified versions of GPL'd software, but \S 1 says that
they can. The ``at no charge'' does not prohibit redistributors from
charging when performing the acts governed by copyright
law,\footnote{Recall that you could by default charge for any acts not
governed by copyright law, because the license controls are confined
by copyright.} but rather that they cannot charge a fee for the
\emph{license itself}. In other words, redistributors of (modified
and unmodified) GPL'd works may charge any amount they choose for
performing the modifications on contract or the act of transferring
the copy to the customer, but they may not charge a separate licensing
fee for the software.

\S 2(b) further states that the software must ``be licensed...to all
third parties.''  This too has led to some confusions, and feeds the
misconception mentioned earlier --- that all modified versions must made
available to the public at large. However, the text here does not say
that. Instead, it says that the licensing under terms of the GPL must
extend to anyone who might, through the distribution chain, receive a copy
of the software. Distribution to all third parties is not mandated here,
but \S 2(b) does require redistributors to license the derivative works in
a way that extends to all third parties who may ultimately receive a
copy of the software.

In summary, \S 2(b) says what terms under which the third parties must
receive this no-charge license. Namely, they receive it ``under the terms
of this License,'' the GPL. When an entity \emph{chooses} to redistribute
a derivative work of GPL'd software, the license of that whole derivative
work must be GPL and only GPL\@. In this manner, \S 2(b) dovetails nicely
with \S 6 (as discussed in Section~\ref{GPLs6} of this tutorial).

\medskip

The final paragraph of \S 2 is worth special mention. It is possible and
quite common to aggregate various software programs together on one
distribution medium. Computer manufacturers do this when they ship a
pre-installed hard drive, and GNU/Linux distribution vendors do this to
give a one-stop CD or URL for a complete operating system with necessary
applications. The GPL very clearly permits such ``mere aggregation'' with
programs under any license. Despite what you hear from its critics, the
GPL is nothing like a virus, not only because the GPL is good for you and
a virus is bad for you, but also because simple contact with a GPL'd
code-base does not impact the license of other programs. Actual effort
must be expended by a programmer to cause a work to fall under the terms
of the GPL. Redistributors are always welcome to simply ship GPL'd
software alongside proprietary software or other unrelated Free Software,
as long as the terms of GPL are adhered to for those packages that are
truly GPL'd.

\section{GPL \S 3: Producing Binaries}
\label{GPL-Section-3}
% FIXME: need name of a novelist who writes very obscurely and obliquely.

Software is a strange beast when compared to other copyrightable works.
It is currently impossible to make a film or a book that can be truly
obscured. Ultimately, the full text of a novel, even one written by
William Faulkner, must presented to the reader as words in some
human-readable language so that they can enjoy the work. A film, even one
directed by David Lynch, must be perceptible by human eyes and ears to
have any value.

Software is not so. While the source code, the human-readable
representation of software is of keen interest to programmers, users and
programmers alike cannot make the proper use of software in that
human-readable form. Binary code --- the ones and zeros that the computer
can understand --- must be predicable and attainable for the software to
be fully useful. Without the binaries, be they in object or executable
form, the software serves only the didactic purposes of computer science.

Under copyright law, binary representations of the software are simply
derivative works of the source code. Applying a systematic process (i.e.,
``compilation'') to a work of source code yields binary code. The binary
code is now a new work of expression fixed in the tangible medium of
electronic file storage.

Therefore, for GPL'd software to be useful, the GPL, since it governs the
rules for creation of derivative works, must grant permission for the
generation of binaries. Furthermore, notwithstanding the relative
popularity of source-based GNU/Linux distributions like Gentoo, users find
it extremely convenient to receive distribution of binary software. Such
distribution is the redistribution of derivative works of the software's
source code. \S 3 addresses the matter of creation and distribution of
binary versions.

Under \S 3, binary versions may be created and distributed under the
terms of \S\S 1--2, so all the material previously discussed applies
here. However, \S 3 must go a bit further. Access to the software's
source code is an incontestable prerequisite for the exercise of the
fundamental freedoms to modify and improve the software. Making even
the most trivial changes to a software program at the binary level is
effectively impossible. \S 3 must ensure that the binaries are never
distributed without the source code, so that these freedoms are passed
through the distribution chain.

\S 3 permits distribution of binaries, and then offers three options for
distribution of source code along with binaries. The most common and the
least complicated is the option given under \S 3(a).

\S 3(a) offers the option to directly accompany the source code alongside
the distribution of the binaries. This is by far the most convenient
option for most distributors, because it means that the source-code
provision obligations are fully completed at the time of binary
distribution (more on that later).

Under \S 3(a), the source code provided must be the ``corresponding source
code.''  Here ``corresponding'' primarily means that the source code
provided must be that code used to produce the binaries being distributed.
That source code must also be ``complete.''  A later paragraph of \S 3
explains in detail what is meant by ``complete.''  In essence, it is all
the material that a programmer of average skill would need to actually use
the source code to produce the binaries she has received. Complete source
is required so that, if the licensee chooses, she should be able to
exercise her freedoms to modify and redistribute changes. Without the
complete source, it would not be possible to make changes that were
actually directly derived from the version received.

Furthermore, \S 3 is defending against a tactic that has in fact been
seen in FSF's GPL enforcement. Under GPL, if you pay a high price for
a copy of GPL'd binaries (which comes with corresponding source, of
course), you have the freedom to redistribute that work at any fee you
choose, or not at all. Sometimes, companies attempt a GPL-violating
cozenage whereby they produce very specialized binaries (perhaps for
an obscure architecture). They then give source code that does
correspond, but withhold the ``incantations'' and build plans they
used to make that source compile into the specialized binaries.
Therefore, \S 3 requires that the source code include ``meta-material'' like
scripts, interface definitions, and other material that is used to
``control compilation and installation'' of the binaries. In this
manner, those further down the distribution chain are assured that
they have the unabated freedom to build their own derivative works
from the sources provided.

FSF (as authors of GPL) realizes that software distribution comes in many
forms. Embedded manufacturers, for example, have the freedom to put
GPL'd software into their PDAs with very tight memory and space
constraints. In such cases, putting the source right alongside the
binaries on the machine itself might not be an option. While it is
recommended that this be the default way that people comply with GPL, the
GPL does provide options when such distribution is infeasible.

\S 3, therefore, allows source code to be provided on any physical
``medium customarily used for software interchange.''  By design, this
phrase covers a broad spectrum. At best, FSF can viably release a new GPL
every ten years or so. Thus, phrases like this must be adaptive to
changes in the technology. When GPL version 2 was first published in June
1991, distribution on magnetic tape was still common, and CD was
relatively new. Today, CD is the default, and for larger systems DVD-R is
gaining adoption. This language must adapt with changing technology.

Meanwhile, the binding created by the word ``customarily'' is key. Many
incorrectly believe that distributing binary on CD and source on the
Internet is acceptable. In the corporate world, it is indeed customary to
simply download CDs worth of data over a T1 or email large file
attachments. However, even today in the USA, many computer users with
CD-ROM drives are not connected to the Internet, and most people connected
to the Internet are connected via a 56K dial-up connection. Downloading
CDs full of data is not customary for them in the least. In some cities
in Africa, computers are becoming more common, but Internet connectivity
is still available only at a few centralized locations. Thus, the
``customs'' here must be normalized for a worldwide userbase. Simply
providing source on the Internet --- while it is a kind, friendly and
useful thing to do --- is not usually sufficient.

Note, however, a major exception to this rule, given by the last paragraph
of \S 3. \emph{If} distribution of the binary files is made only on the
Internet (i.e., ``from a designated place''), \emph{then} simply providing
the source code right alongside the binaries in the same place is
sufficient to comply with \S 3.

\medskip

As is shown above, Under \S 3(a), embedded manufacturers can put the
binaries on the device and ship the source code along on a CD\@. However,
sometimes this turns out to be too costly. Including a CD with every
device could prove too costly, and may practically (although not legally)
prohibit using GPL'd software. For this situation and others like it, \S
3(b) is available.

\S 3(b) allows a distributor of binaries to instead provide a written
offer for source code alongside those binaries. This is useful in two
specific ways. First, it may turn out that most users do not request the
source, and thus the cost of producing the CDs is saved --- a financial
and environmental windfall. In addition, along with a \S 3(b) compliant
offer for source, a binary distributor might choose to \emph{also} give a
URL for source code. Many who would otherwise need a CD with source might
turn out to have those coveted high bandwidth connections, and are able to
download the source instead --- again yielding environmental and financial
windfalls.

However, note that regardless of how many users prefer to get the
source online, \S 3(b) does place lasting long-term obligations on the
binary distributor. The binary distributor must be prepared to honor
that offer for source for three years and ship it out (just as they
would have had to do under \S 3(a)) at a moment's notice when they
receive such a request. There is real organizational cost here:
support engineers must be trained how to route source requests, and
source CD images for every release version for the last three years
must be kept on hand to burn such CDs quickly. The requests might not
even come from actual customers; the offer for source must be valid
for ``any third party.''

That phrase is another place where some get confused --- thinking again
that full public distribution of source is required. The offer for source
must be valid for ``any third party'' because of the freedoms of
redistribution granted by \S\S 1--2. A company may ship a binary image
and an offer for source to only one customer. However, under GPL, that
customer has the right to redistribute that software to the world if she
likes. When she does, that customer has an obligation to make sure that
those who receive the software from her can exercise their freedoms under
GPL --- including the freedom to modify, rebuild, and redistribute the
source code.

\S 3(c) is created to save her some trouble, because by itself \S 3(b)
would unfairly favor large companies. \S 3(b) allows the
separation of the binary software from the key tool that people can use
to exercise their freedom. The GPL permits this separation because it is
good for redistributors, and those users who turn out not to need the
source. However, to ensure equal rights for all software users, anyone
along the distribution chain must have the right to get the source and
exercise those freedoms that require it.

Meanwhile, \S 3(b)'s compromise primarily benefits companies who
distribute binary software commercially. Without \S 3(c), that benefit
would be at the detriment of the companies' customers; the burden of
source code provision would be unfairly shifted to the companies'
customers. A customer, who had received binaries with a \S 3(b)-compliant
offer, would be required under GPL (sans \S 3(c)) to acquire the source,
merely to give a copy of the software to a friend who needed it. \S 3(c)
reshifts this burden to entity who benefits from \S 3(b).

\S 3(c) allows those who undertake \emph{noncommercial} distribution to
simply pass along a \S 3(b)-compliant source code offer. The customer who
wishes to give a copy to her friend can now do so without provisioning the
source, as long as she gives that offer to her friend. By contrast, if
she wanted to go into business for herself selling CDs of that software,
she would have to acquire the source and either comply via \S 3(a), or
write her own \S 3(b)-compliant source offer.

This process is precisely the reason why a \S 3(b) source offer must be
valid for all third parties. At the time the offer is made, there is no
way of knowing who might end up noncommercially receiving a copy of the
software. Companies who choose to comply via \S 3(b) must thus be
prepared to honor all incoming source code requests. For this and the
many other additional necessary complications under \S\S 3(b--c), it is
only rarely a better option than complying via \S 3(a).

%%%%%%%%%%%%%%%%%%%%%%%%%%%%%%%%%%%%%%%%%%%%%%%%%%%%%%%%%%%%%%%%%%%%%%%%%%%%%%%
\chapter{The Implied Patent Grant in GPL}

We digress again briefly from our section-by-section consideration of GPL
to consider the interaction between the terms of GPL and patent law. The
GPL, despite being silent with respect to patents, actually confers on its
licensees more rights to a licensor's patents than those licenses that
purport to address the issue. This is the case because patent law, under
the doctrine of implied license, gives to each distributee of a patented
article a license from the distributor to practice any patent claims owned
or held by the distributor that cover the distributed article. The
implied license also extends to any patent claims owned or held by the
distributor that cover ``reasonably contemplated uses'' of the patented
article. To quote the Federal Circuit Court of Appeals, the highest court
for patent cases other than the Supreme Court:

\begin{quotation}
Generally, when a seller sells a product without restriction, it in
effect promises the purchaser that in exchange for the price paid, it will
not interfere with the purchaser's full enjoyment of the product
purchased. The buyer has an implied license under any patents of the
seller that dominate the product or any uses of the product to which the
parties might reasonably contemplate the product will be put.
\end{quotation}
Hewlett-Packard Co. v. Repeat-O-Type Stencil Mfg. Corp., Inc., 123 F.3d
1445 (Fed. Cir. 1997).

Of course, Free Software is licensed, not sold, and there are indeed
restrictions placed on the licensee, but those differences are not likely
to prevent the application of the implied license doctrine to Free
Software, because software licensed under the GPL grants the licensee the
right to make, use, and sell the software, each of which are exclusive
rights of a patent holder. Therefore, although the GPL does not expressly
grant the licensee the right to do those things under any patents the
licensor may have that cover the software or its reasonably contemplated
uses, by licensing the software under the GPL, the distributor impliedly
licenses those patents to the GPL licensee with respect to the GPL
licensed software.

An interesting issue regarding this implied patent license of GPL'd
software is what would be considered ``uses of the [software] to which
the parties might reasonably contemplate the product will be put.'' A
clever advocate may argue that the implied license granted by GPL is
larger in scope than the express license in other Free Software
licenses with express patent grants, in that, the patent license
clause of many of those licenses are specifically limited to the
patent claims covered by the code as licensed by the patentee.

To the contrary, GPL's implied patent license grants the GPL licensee a
patent license to do much more than just that because the GPL licensee,
under the doctrine of implied patent license, is free to practice any
patent claims held by the licensor that cover ``reasonably contemplated
uses'' of the GPL'd code, which may very well include creation and
distribution of derivative works since the GPL's terms, under which the
patented code is distributed, expressly permits such activity.

Further supporting this result is the Federal Circuit's pronouncement that
the recipient of a patented article has, not only an implied license to
make, use, and sell the article, but also an implied patent license to
repair the article to enable it to function properly, Bottom Line Mgmt.,
Inc. v. Pan Man, Inc., 228 F.3d 1352 (Fed. Cir. 2000). Additionally, the
Federal Circuit extended that rule to include any future recipients of the
patented article, not just the direct recipient from the distributor.
This theory comports well with the idea of Free Software, whereby software
is distributed amongst many entities within the community for the purpose
of constant evolution and improvement. In this way, the law of implied
patent license used by the GPL ensures that the community mutually
benefits from the licensing of patents to any single community member.

Note that simply because GPL'd software has an implied patent license does
not mean that any patents held by a distributor of GPL'd code become
worthless. To the contrary, the patents are still valid and enforceable
against either:

\begin{enumerate}
 \renewcommand{\theenumi}{\alph{enumi}}
 \renewcommand{\labelenumi}{\textup{(\theenumi)}}

\item any software other than that licensed under the GPL by the patent
  holder, and

\item any party that does not comply with the GPL
with respect to the licensed software.
\end{enumerate}

\newcommand{\compB}{$\mathcal{B}$}
\newcommand{\compA}{$\mathcal{A}$}

For example, if Company \compA{} has a patent on advanced Web browsing, but
also licenses a Web browsing software program under the GPL, then it
cannot assert the patent against any party that takes a license to its
program under the GPL. However, if a party uses that program without
complying with the GPL, then Company \compA{} can assert, not just copyright
infringement claims against the non-GPL-compliant party, but also
infringement of the patent, because the implied patent license only
extends to use of the software in accordance with the GPL. Further, if
Company \compB{} distributes a competitive advanced Web browsing program,
Company \compA{} is free to assert its patent against any user or
distributor of that product. It is irrelevant whether Company \compB's
program is distributed under the GPL, as Company \compB{} can not grant
implied licenses to Company \compA's patent.

This result also reassures companies that they need not fear losing their
proprietary value in patents to competitors through the GPL implied patent
license, as only those competitors who adopt and comply with the GPL's
terms can benefit from the implied patent license. To continue the
example above, Company \compB{} does not receive a free ride on Company
\compA's patent, as Company \compB{} has not licensed-in and then
redistributed Company A's advanced Web browser under the GPL. If Company
\compB{} does do that, however, Company \compA{} still has not lost
competitive advantage against Company \compB{}, as Company \compB{} must then,
when it re-distributes Company \compA's program, grant an implied license
to any of its patents that cover the program. Further, if Company \compB{}
relicenses an improved version of Company A's program, it must do so under
the GPL, meaning that any patents it holds that cover the improved version
are impliedly licensed to any licensee. As such, the only way Company
\compB{} can benefit from Company \compA's implied patent license, is if it,
itself, distributes Company \compA's software program and grants an
implied patent license to any of its patents that cover that program.

%%%%%%%%%%%%%%%%%%%%%%%%%%%%%%%%%%%%%%%%%%%%%%%%%%%%%%%%%%%%%%%%%%%%%%%%%%%%%%%
\chapter{Defending Freedom on Many Fronts}

Chapters~\ref{run-and-verbatim} and ~\ref{source-and-binary} presented the
core freedom-defending provisions of GPL\@, which are in \S\S 0--3. \S\S
4--7 of the GPL are designed to ensure that \S\S 0--3 are not infringed,
are enforceable, are kept to the confines of copyright law, and are not
trumped by other copyright agreements or components of other entirely
separate legal systems. In short, while \S\S 0--3 are the parts of the
license that defend the freedoms of users and programmers, \S\S 4--7 are
the parts of the license that keep the playing field clear so that \S\S
0--3 can do their jobs.

\section{GPL \S 4: Termination on Violation}
\label{GPLs4}

\S 4 is GPL's termination clause. Upon first examination, it seems
strange that a license with the goal of defending users' and programmers'
freedoms for perpetuity in an irrevocable way would have such a clause.
However, upon further examination, the difference between irrevocability
and this termination clause becomes clear.

The GPL is irrevocable in the sense that once a copyright holder grants
rights for someone to copy, modify and redistribute the software under
terms of the GPL, they cannot later revoke that grant. Since the GPL has
no provision allowing the copyright holder to take such a prerogative, the
license is granted as long as the copyright remains in effect.\footnote{In
  the USA, due to unfortunate legislation, the length of copyright is
  nearly perpetual, even though the Constitution forbids perpetual
  copyright.} The copyright holder has the right to relicense the same
work under different licenses (see Section~\ref{Proprietary Relicensing}
of this tutorial), or to stop distributing the GPL'd version (assuming \S
3(b) was never used), but she may not revoke the rights under GPL
already granted.

In fact, when an entity looses their right to copy, modify and distribute
GPL'd software, it is because of their \emph{own actions}, not that of
the copyright holder. The copyright holder does not decided when \S 4
termination occurs (if ever), the actions of the licensee does.

Under copyright law, the GPL has granted various rights and freedoms to
the licensee to perform specific types of copying, modification, and
redistribution. By default, all other types of copying, modification, and
redistribution are prohibited. \S 4 says that if you undertake any of
those other types (e.g., redistributing binary-only in violation of \S 3),
then all rights under the license --- even those otherwise permitted for
those who have not violated --- terminate automatically.

\S 4 gives GPL teeth. If licensees fail to adhere to the license, then
they are stuck. They must completely cease and desist from all
copying, modification and distribution of that GPL'd software.

At that point, violating licensees must gain the forgiveness of the
copyright holder to have their rights restored. Alternatively, they could
negotiate another agreement, separate from GPL, with the copyright
holder. Both are common practice.

At FSF, it is part of the mission to spread software freedom. When FSF
enforces GPL, the goal is to bring the violator back into compliance as
quickly as possible, and redress the damage caused by the violation.
That is FSF's steadfast position in a violation negotiation --- comply
with the license and respect freedom.

However, other entities who do not share the full ethos of software
freedom as institutionalized by FSF pursue GPL violations differently.
MySQL AB, a company that produces the GPL'd MySQL database, upon
discovering GPL violations typically negotiates a proprietary software
license separately for a fee. While this practice is not one that FSF
would ever consider undertaking or even endorsing, it is a legal way for
copyright holders to proceed.

\section{GPL \S 5: Acceptance, Copyright Style}
\label{GPLs5}

\S 5 brings us to perhaps the most fundamental misconception and common
confusion about GPL\@. Because of the prevalence of proprietary software,
most users, programmers, and lawyers alike tend to be more familiar with
EULAs. EULAs are believed by their authors to be contracts, requiring
formal agreement between the licensee and the software distributor to be
valid. This has led to mechanisms like ``shrink-wrap'' and ``click-wrap''
as mechanisms to perform acceptance ceremonies with EULAs.

The GPL does not need contract law to ``transfer rights.''  No rights
are transfered between parties. By contrast, the GPL is a permission
slip to undertake activities that would otherwise have been prohibited
by copyright law. As such, it needs no acceptance ceremony; the
licensee is not even required to accept the license.

However, without the GPL, the activities of copying, modifying and
distributing the software would have otherwise been prohibited. So, the
GPL says that you only accepted the license by undertaking activities that
you would have otherwise been prohibited without your license under GPL\@.
This is a certainly subtle point, and requires a mindset quite different
from the contractual approach taken by EULA authors.

An interesting side benefit to \S 5 is that the bulk of users of Free
Software are not required to accept the license. Undertaking fair and
unregulated use of the work, for example, does not bind you to the GPL,
since you are not engaging in activity that is otherwise controlled by
copyright law. Only when you engage in those activities that might have an
impact on the freedom of others does license acceptance occur, and the
terms begin to bind you to fair and equitable sharing of the software. In
other words, the GPL only kicks in when it needs to for the sake of
freedom.

\section{GPL \S 6: GPL, My One and Only}
\label{GPLs6}

A point that was glossed over in Section~\ref{GPLs4}'s discussion of \S 4
was the irrevocable nature of the GPL\@. The GPL is indeed irrevocable,
and it is made so formally by \S 6.

The first sentence in \S 6 ensures that as software propagates down the
distribution chain, that each licensor can pass along the license to each
new licensee. Under \S 6, the act of distributing automatically grants a
license from the original licensor to the next recipient. This creates a
chain of grants that ensure that everyone in the distribution has rights
under the GPL\@. In a mathematical sense, this bounds the bottom ---
making sure that future licensees get no fewer rights than the licensee before.

The second sentence of \S 6 does the opposite; it bounds from the top. It
prohibits any licensor along the distribution chain from placing
additional restrictions on the user. In other words, no additional
requirements may trump the rights and freedoms given by GPL\@.

The final sentence of \S 6 makes it abundantly clear that no individual
entity in the distribution chain is responsible for the compliance of any
other. This is particularly important for noncommercial users who have
passed along a source offer under \S 3(c), as they cannot be assured that
the issuer of the offer will honor their \S 3 obligations.

In short, \S 6 says that your license for the software is your one and
only copyright license allowing you to copy, modify and distribute the
software.

\section{GPL \S 7: ``Give Software Liberty or Give It Death!''}
\label{GPLs7}

In essence, \S 7 is a verbosely worded way of saying for non-copyright
systems what \S 6 says for copyright. If there exists any reason that a
distributor knows of that would prohibit later licensees from exercising
their full rights under GPL, then distribution is prohibited.

Originally, this was designed as the title of this section suggests --- as
a last ditch effort to make sure that freedom was upheld. However, in
modern times, it has come to give much more. Now that the body of GPL'd
software is so large, patent holders who would want to be distributors of
GPL'd software have a tough choice. They must choose between avoiding
distribution of GPL'd software that exercises the teachings of their
patents, or grant a royalty-free, irrevocable, non-exclusive license to
those patents. Many companies, including IBM, the largest patent holder
in the world, have chosen the latter.

Thus, \S 7 rarely gives software death by stopping its distribution.
Instead, it is inspiring patent holders to share their patents in the same
freedom-defending way that they share their copyrighted works.

\section{GPL \S 8: Excluding Unfreedonia}
\label{GPLs8}

\S 8 is rarely used by copyright holders. Its intention is that if a
particular country, say Unfreedonia, grants particular patents or allows
copyrighted interfaces (no country to our knowledge even permits those
yet), that the GPL'd software can continue in free and unabated
distribution in the countries where such controls do not exist.

It is a partial ``out'' from \S 7. Without \S 8, if a copyright holder
knew of a patent in a particular country licensed in a GPL-incompatible
way, then she could not distribute under GPL, because the work could
legitimately end up in the hands of citizens of Unfreedonia.

It is an inevitable but sad reality that some countries are freer than
others. \S 8 exists to permit distribution in those countries that are
free without otherwise negating parts of the license.

%%%%%%%%%%%%%%%%%%%%%%%%%%%%%%%%%%%%%%%%%%%%%%%%%%%%%%%%%%%%%%%%%%%%%%%%%%%%%%%
\chapter{Odds, Ends, and Absolutely No Warranty}

\S 0--7 constitute the freedom-defending terms of the GPL. The remainder
of the GPL handles administrivia and issues concerning warranties and
liability.

\section{GPL \S 9: FSF as Stewards of GPL}
\label{GPLs9}

FSF reserves the exclusive right to publish future versions of the GPL\@;
\S 9 expresses this. While the stewardship of the copyrights on the body
of GPL'd software around the world is shared among thousands of
individuals and organizations, the license itself needs a single steward.
Forking of the code is often regrettable but basically innocuous. Forking
of licensing is disastrous.

FSF has only released two versions of GPL --- in 1989 and 1991. GPL
version 3 is under current internal drafting. FSF's plan is to have a
long and engaging comment period. The goal of GPL is to defend freedom, and
a gigantic community depends on that freedom now. FSF hopes to take all
stakeholders' opinions under advisement.

\section{GPL \S 10: Relicensing Permitted}
\label{GPLs10}

\S 10 reminds the licensee of what is already implied by the nature of
copyright law. Namely, the copyright holder of a particular software
program has the prerogative to grant alternative agreements under separate
copyright licenses.

\section{GPL \S 11: No Warranty}
\label{GPLs11}

All warranty disclaimer language tends to be shouted in all capital
letters. Apparently, there was once a case where the disclaimer language
of an agreement was negated because it was not ``conspicuous'' to one of
the parties. Therefore, to make such language ``conspicuous,'' people
started placing it in bold or capitalizing the entire text. It now seems
to be voodoo tradition of warranty disclaimer writing.

Some have argued the GPL is unenforceable in some jurisdictions because
its disclaimer of warranties is impermissibly broad. However, \S 11
contains a jurisdictional savings provision, which states that it is to be
interpreted only as broadly as allowed by applicable law. Such a
provision ensures that both it, and the entire GPL, is enforceable in any
jurisdiction, regardless of any particular law regarding the
permissibility of certain warranty disclaimers.

Finally, one important point to remember when reading \S 11 is that \S 1
permits the sale of warranty as an additional service, which \S 11 affirms.

\section{GPL, \S 12: Limitation of Liability}
\label{GPLs12}

There are many types of warranties, and in some jurisdictions some of them
cannot be disclaimed. Therefore, usually agreements will have both a
warranty disclaimer and a limitation of liability, as we have in \S 12. \S
11 thus gets rid of all implied warranties that can legally be
disavowed. \S 12, in turn, limits the liability of the actor for any
warranties that cannot legally be disclaimed in a particular jurisdiction.

Again, some have argued the GPL is unenforceable in some jurisdictions
because its limitation of liability is impermissibly broad. However, \S
12, just like its sister, \S 11, contains a jurisdictional savings
provision, which states that it is to be interpreted only as broadly as
allowed by applicable law. As stated above, such a provision ensures that
both \S 12, and the entire GPL, is enforceable in any jurisdiction,
regardless of any particular law regarding the permissibility of limiting
liability.

So end the terms and conditions of the GNU General Public License.

%%%%%%%%%%%%%%%%%%%%%%%%%%%%%%%%%%%%%%%%%%%%%%%%%%%%%%%%%%%%%%%%%%%%%%%%%%%%%%%
\chapter{The Lesser GPL}

As we have seen in our consideration of the GPL, its text is specifically
designed to cover all possible derivative works under copyright law. Our
goal in designing GPL was to make sure that any derivative work of GPL'd
software was itself released under GPL when distributed. Reaching as far
as copyright law will allow is the most direct way to reach that goal.

However, while the strategic goal is to bring as much Free Software
into the world as possible, particular tactical considerations
regarding software freedom dictate different means. Extending the
copyleft effect as far as copyright law allows is not always the most
prudent course in reaching the goal. In particular situations, even
those of us with the goal of building a world where all published
software is Free Software realize that full copyleft does not best
serve us. The GNU Lesser General Public License (``GNU LGPL'') was
designed as a solution for such situations.

\section{The First LGPL'd Program}

The first example that FSF encountered where such altered tactics were
needed was when work began on the GNU C Library. The GNU C Library would
become (and today, now is) a drop-in replacement for existing C libraries.
On a Unix-like operating system, C is the lingua franca and the C library
is an essential component for all programs. It is extremely difficult to
construct a program that will run with ease on a Unix-like operating
system without making use of services provided by the C library --- even
if the program is written in a language other than C\@. Effectively, all
user application programs that run on any modern Unix-like system must
make use of the C library.

By the time work began on the GNU implementation of the C libraries, there
were already many C libraries in existence from a variety of vendors.
Every proprietary Unix vendor had one, and many third parties produced
smaller versions for special purpose use. However, our goal was to create
a C library that would provide equivalent functionality to these other C
libraries on a Free Software operating system (which in fact happens today
on modern GNU/Linux systems, which all use the GNU C Library).

Unlike existing GNU application software, however, the licensing
implications of releasing the GNU C Library (``glibc'') under GPL were
somewhat different. Applications released under GPL would never
themselves become part of proprietary software. However, if glibc were
released under GPL, it would require that any application distributed for
the GNU/Linux platform be released under GPL\@.

Since all applications on a Unix-like system depend on the C library, it
means that they must link with that library to function on the system. In
other words, all applications running on a Unix-like system must be
combined with the C library to form a new whole derivative work that is
composed of the original application and the C library. Thus, if glibc
were GPL'd, each and every application distributed for use on GNU/Linux
would also need to be GPL'd, since to even function, such applications
would need to be combined into larger derivative works by linking with
glibc.

At first glance, such an outcome seems like a windfall for Free Software
advocates, since it stops all proprietary software development on
GNU/Linux systems. However, the outcome is a bit more subtle. In a world
where many C libraries already exist, many of which could easily be ported
to GNU/Linux, a GPL'd glibc would be unlikely to succeed. Proprietary
vendors would see the excellent opportunity to license their C libraries
to anyone who wished to write proprietary software for GNU/Linux systems.
The de-facto standard for the C library on GNU/Linux would likely be not
glibc, but the most popular proprietary one.

Meanwhile, the actual goal of releasing glibc under GPL --- to ensure no
proprietary applications on GNU/Linux --- would be unattainable in this
scenario. Furthermore, users of those proprietary applications would also
be users of a proprietary C library, not the Free glibc.

The Lesser GPL was initially conceived to handle this scenario. It was
clear that the existence of proprietary applications for GNU/Linux was
inevitable. Since there were so many C libraries already in existence, a
new one under GPL would not stop that tide. However, if the new C library
were released under a license that permitted proprietary applications
to link with it, but made sure that the library itself remained Free,
an ancillary goal could be met. Users of proprietary applications, while
they would not have the freedom to copy, share, modify and redistribute
the application itself, would have the freedom to do so with respect to
the C library.

There was no way the license of glibc could stop or even slow the creation
of proprietary applications on GNU/Linux. However, loosening the
restrictions on the licensing of glibc ensured that nearly all proprietary
applications at least used a Free C library rather than a proprietary one.
This trade-off is central to the reasoning behind the LGPL\@.

Of course, many people who use the LGPL today are not thinking in these
terms. In fact, they are often choosing the LGPL because they are looking
for a ``compromise'' between the GPL and the X11-style liberal licensing.
However, understanding FSF's reasoning behind the creation of the LGPL is
helpful when studying the license.


\section{What's the Same?}

Much of the text of the LGPL is identical to the GPL\@. As we begin our
discussion of the LGPL, we will first eliminate the sections that are
identical, or that have the minor modification changing the word
``Program'' to ``Library.''

First, \S 1 of LGPL, the rules for verbatim copying of source, are
equivalent to those in GPL's \S 1.

Second, \S 8 of LGPL is equivalent \S 4 of GPL\@. In both licenses, this
section handles termination in precisely the same manner.

\S 9 in LGPL is equivalent to \S 5 in GPL\@. Both sections assert that
the license is a copyright license, and handle the acceptance of those
copyright terms.

LGPL's \S 10 is equivalent to GPL's \S 6. They both protect the
distribution system of Free Software under these licenses, to ensure that
up, down, and throughout the distribution chain, each recipient of the
software receives identical rights under the license and no other
restrictions are imposed.

LGPL's \S 11 is GPL's \S 7. As discussed, it is used to ensure that
other claims and legal realities, such as patent licenses and court
judgments, do not trump the rights and permissions granted by these
licenses, and requires that distribution be halted if such a trump is
known to exist.

LGPL's \S 12 adds the same features as GPL's \S 8. These sections are
used to allow original copyright holders to forbid distribution in
countries with draconian laws that would otherwise contradict these
licenses.

LGPL's \S 13 sets up FSF as the steward of the LGPL, just as GPL's \S 9
does for GPL. Meanwhile, LGPL's \S 14 reminds licensees that copyright
holders can grant exceptions to the terms of LGPL, just as GPL's \S 10
reminds licensees of the same thing.

Finally, the assertions of no warranty and limitations of liability are
identical; thus LGPL's \S 15 and \S 16 are the same as GPL's \S 11 and \S
12.

As we see, the entire latter half of the license is identical.
The parts which set up the legal boundaries and meta-rules for the license
are the same. It is our intent that the two licenses operate under the
same legal mechanisms and are enforced precisely the same way.

We strike a difference only in the early portions of the license.
Namely, in the LGPL we go into deeper detail of granting various permissions to
create derivative works, so the redistributors can make
some proprietary derivatives. Since we simply do not allow the
license to stretch as far as copyright law does regarding what
derivative works must be relicensed under the same terms, we must go
further to explain which derivative works we will allow to be
proprietary. Thus, we'll see that the front matter of the LGPL is a
bit more wordy and detailed with regards to the permissions granted to
those who modify or redistribute the software.

\section{Additions to the Preamble}

Most of LGPL's Preamble is identical, but the last seven paragraphs
introduce the concepts and reasoning behind creation of the license,
presenting a more generalized and briefer version of the story with which
we began our consideration of LGPL\@.

In short, FSF designed LGPL for those edge cases where the freedom of the
public can better be served by a more lax licensing system. FSF doesn't
encourage use of LGPL automatically for any software that happens to be a
library; rather, FSF suggests that it only be used in specific cases, such
as the following:

\begin{itemize}

\item To encourage the widest possible use of a Free Software library, so
  it becomes a de-facto standard over similar, although not
  interface-identical, proprietary alternatives

\item To encourage use of a Free Software library that already has
  interface-identical proprietary competitors that are more developed

\item To allow a greater number of users to get freedom, by encouraging
  proprietary companies to pick a Free alternative for its otherwise
  proprietary products

\end{itemize}

LGPL's preamble sets forth the limits to which the license seeks to go in
chasing these goals. LGPL is designed to ensure that users who happen to
acquire software linked with such libraries have full freedoms with
respect to that library. They should have the ability to upgrade to a newer
or modified Free version or to make their own modifications, even if they
cannot modify the primary software program that links to that library.

Finally, the preamble introduces two terms used throughout the license to
clarify between the different types of derivative works: ``works that use
the library,'' and ``works based on the library.''  Unlike GPL, LGPL must
draw some lines regarding derivative works. We do this here in this
license because we specifically seek to liberalize the rights afforded to
those who make derivative works. In GPL, we reach as far as copyright law
allows. In LGPL, we want to draw a line that allows some derivative works
copyright law would otherwise prohibit if the copyright holder exercised
his full permitted controls over the work.

\section{A Work that Uses the Library}

In the effort to allow certain proprietary derivative works and prohibit
others, LGPL distinguishes between two classes of derivative works:
``works based on the library,'' and ``works that use the library.''  The
distinction is drawn on the bright line of binary (or runtime) derivative
works and source code derivatives. We will first consider the definition
of a ``work that uses the library,'' which is set forth in LGPL \S 5.

We noted in our discussion of GPL \S 3 (discussed in
Section~\ref{GPL-Section-3} of this document) that binary programs when
compiled and linked with GPL'd software are derivative works of that GPL'd
software. This includes both linking that happens at compile-time (when
the binary is created) or at runtime (when the binary -- including library
and main program both -- is loaded into memory by the user). In GPL,
binary derivative works are controlled by the terms of the license (in GPL
\S 3), and distributors of such binary derivatives must release full
corresponding source\@.

In the case of LGPL, these are precisely the types of derivative works
we wish to permit. This scenario, defined in LGPL as ``a work that uses
the library,'' works as follows:

\newcommand{\workl}{$\mathcal{L}$}
\newcommand{\lplusi}{$\mathcal{L\!\!+\!\!I}$}

\begin{itemize}

\item A new copyright holder creates a separate and independent work,
  \worki{}, that makes interface calls (e.g., function calls) to the
  LGPL'd work, called \workl{}, whose copyright is held by some other
  party. Note that since \worki{} and \workl{} are separate and
  independent works, there is no copyright obligation on this new copyright
  holder with regard to the licensing of \worki{}, at least with regard to
  the source code.

\item The new copyright holder, for her software to be useful, realizes
  that it cannot run without combining \worki{} and \workl{}.
  Specifically, when she creates a running binary program, that running
  binary must be a derivative work, called \lplusi{}, that the user can
  run.

\item Since \lplusi{} is a derivative work of both \worki{} and \workl{},
  the license of \workl{} (the LGPL) can put restrictions on the license
  of \lplusi{}. In fact, this is what LGPL does.

\end{itemize}

We will talk about the specific restrictions LGPL places on ``works
that use the library'' in detail in Section~\ref{lgpl-section-6}. For
now, focus on the logic related to how the LGPL places requirements on
the license of \lplusi{}. Note, first of all, the similarity between
this explanation and that in Section~\ref{separate-and-independent},
which discussed the combination of otherwise separate and independent
works with GPL'd code. Effectively, what LGPL does is say that when a
new work is otherwise separate and independent, but has interface
calls out to an LGPL'd library, then it is considered a ``work that
uses the library.''

In addition, the only reason that LGPL has any control over the licensing
of a ``work that uses the library'' is for the same reason that GPL has
some say over separate and independent works. Namely, such controls exist
because the {\em binary combination\/} (\lplusi{}) that must be created to
make the separate work (\worki{}) at all useful is a derivative work of
the LGPL'd software (\workl{}).

Thus, a two-question test that will help indicate if a particular work is
a ``work that uses the library'' under LGPL is as follows:

\begin{enumerate}

\item Is the source code of the new copyrighted work, \worki{}, a
  completely independent work that stands by itself, and includes no
  source code from \workl{}?

\item When the source code is compiled, does it create a derivative work
  by combining with \workl{}, either by static (compile-time) or dynamic
  (runtime) linking, to create a new binary work, \lplusi{}?
\end{enumerate}

If the answers to both questions are ``yes,'' then \worki{} is most likely
a ``work that uses the library.''  If the answer to the first question
``yes,'' but the answer to the second question is ``no,'' then most likely
\worki{} is neither a ``work that uses the library'' nor a ``work based on
the library.''  If the answer to the first question is ``no,'' but the
answer to the second question is ``yes,'' then an investigation into
whether or not \worki{} is in fact a ``work based on the library'' is
warranted.

\section{A Work Based on the Library}

In short, a ``work based on the library'' could be defined as any
derivative work of LGPL'd software that cannot otherwise fit the
definition of a ``work that uses the library.''  A ``work based on the
library'' extends the full width and depth of copyright derivative works,
in the same sense that GPL does.

Most typically, one creates a ``work based on the library'' by directly
modifying the source of the library. Such a work could also be created by
tightly integrating new software with the library. The lines are no doubt
fuzzy, just as they are with GPL'd works, since copyright law gives us no
litmus test for derivative works of a software program.

Thus, the test to use when considering whether something is a ``work
based on the library'' is as follows:

\begin{enumerate}

\item Is the new work, when in source form, a derivative work under
  copyright law of the LGPL'd work?

\item Is there no way in which the new work fits the definition of a
  ``work that uses the library''?
\end{enumerate}


If the answer is ``yes'' to both these questions, then you most likely
have a ``work based on the library.''  If the answer is ``no'' to the
first but ``yes'' to the second, you are in a gray area between ``work
based on the library'' and a ``work that uses the library.''

In our years of work with the LGPL, however, we have never seen a work
of software that was not clearly one or the other; the line is quite
bright. At times, though, we have seen cases where a derivative work
appeared in some ways to be a work that used the library and in other
ways a work based on the library. We overcame this problem by
dividing the work into smaller subunits. It was soon discovered that
what we actually had were three distinct components: the original
LGPL'd work, a specific set of works that used that library, and a
specific set of works that were based on the library. Once such
distinctions are established, the licensing for each component can be
considered independently and the LGPL applied to each work as
prescribed.


\section{Subtleties in Works that Use the Library}

In our discussion of the definition of ``works that use the library,'' we
left out a few more complex details that relate to lower-level programming
details. The fourth paragraph of LGPL's \S 5 covers these complexities,
and it has been a source of great confusion. Part of the confusion comes
because a deep understanding of how compiler programs work is nearly
mandatory to grasp the subtle nature of what \S 5, \P 4 seeks to
cover. It helps some to note that this is a border case that we cover in
the license only so that when such a border case is hit, the implications
of using LGPL continue in the expected way.

To understand this subtle point, we must recall the way that a compiler
operates. The compiler first generates object code, which are the binary
representations of various programming modules. Each of those modules is
usually not useful by itself; it becomes useful to a user of a full program
when those modules are {\em linked\/} into a full binary executable.

As we have discussed, the assembly of modules can happen at compile-time
or at runtime. Legally, there is no distinction between the two --- both
create a derivative work by copying and combining portions of one work and
mixing them with another. However, under LGPL, there is a case in the
compilation process where the legal implications are different.
Specifically, while we know that a ``work that uses the library'' is one
whose final binary is a derivative work, but whose source is not, there
are cases where the object code --- that intermediate step between source
and final binary --- is a derivative work created by copying verbatim code
from the LGPL'd software.

For efficiency, when a compiler turns source code into object code, it
sometimes places literal portions of the copyrighted library code into the
object code for an otherwise separate independent work. In the normal
scenario, the derivative would not be created until final assembly and
linking of the executable occurred. However, when the compiler does this
efficiency optimization, at the intermediate object code step, a
derivative work is created.

LGPL's \S 5, \P 4 is designed to handle this specific case. The intent of
the license is clearly that simply compiling software to ``make use'' of
the library does not in itself cause the compiled work to be a ``work
based on the library.''  However, since the compiler copies verbatim,
copyrighted portions of the library into the object code for the otherwise
separate and independent work, it would actually cause that object file to be a
``work based on the library.''  It is not FSF's intent that a mere
compilation idiosyncrasy would change the requirements on the users of the
LGPL'd software. This paragraph removes that restriction, allowing the
implications of the license to be the same regardless of the specific
mechanisms the compiler uses underneath to create the ``work that uses the
library.''

As it turns out, we have only once had anyone worry about this specific
idiosyncrasy, because that particular vendor wanted to ship object code
(rather than final binaries) to their customers and was worried about
this edge condition. The intent of clarifying this edge condition is
primarily to quell the worries of software engineers who understand the
level of verbatim code copying that a compiler often does, and to help
them understand that the full implications of LGPL are the same regardless
of the details of the compilation progress.

\section{LGPL \S 6: Distributing Works that Use the Library}
\label{lgpl-section-6}
Now that we have established a good working definition of works that
``use'' and works that ``are based on'' the library, we will consider the
rules for distributing these two different works.

The rules for distributing ``works that use the library'' are covered in
\S 6 of LGPL\@. \S 6 is much like GPL's \S 3, as it requires the release
of source when a binary version of the LGPL'd software is released. Of
course, it only requires that source code for the library itself be made
available. The work that ``uses'' the library need not be provided in
source form. However, there are also conditions in LGPL \S 6 to make sure
that a user who wishes to modify or update the library can do so.

LGPL \S 6 lists five choices with regard to supplying library source
and granting the freedom to modify that library source to users. We
will first consider the option given by \S 6(b), which describes the
most common way currently used for LGPL compliance on a ``work that
uses the library.''

\S 6(b) allows the distributor of a ``work that uses the library'' to
simply use a dynamically linked, shared library mechanism to link with the
library. This is by far the easiest and most straightforward option for
distribution. In this case, the executable of the work that uses the
library will contain only the ``stub code'' that is put in place by the
shared library mechanism, and at runtime the executable will combine with
the shared version of the library already resident on the user's computer.
If such a mechanism is used, it must allow the user to upgrade and
replace the library with interface-compatible versions and still be able
to use the ``work that uses the library.''  However, all modern shared
library mechanisms function as such, and thus \S 6(b) is the simplest
option, since it does not even require that the distributor of the ``work
based on the library'' ship copies of the library itself.

\S 6(a) is the option to use when, for some reason, a shared library
mechanism cannot be used. It requires that the source for the library be
included, in the typical GPL fashion, but it also has a requirement beyond
that. The user must be able to exercise her freedom to modify the library
to its fullest extent, and that means recombining it with the ``work based
on the library.''  If the full binary is linked without a shared library
mechanism, the user must have available the object code for the ``work
based on the library,'' so that the user can relink the application and
build a new binary.

The remaining options in \S 6 are very similar to the other choices
provided by GPL \S 3. There are some additional options, but time does
not permit us in this course to go into those additional options. In
almost all cases of distribution under LGPL, either \S 6(a) or \S 6(b) are
exercised.

\section{Distribution of Works Based on the Library}

Essentially, ``works based on the library'' must be distributed under the
same conditions as works under full GPL\@. In fact, we note that LGPL's
\S 2 is nearly identical in its terms and requirements to GPL's \S 2.
There are again subtle differences and additions, which time does not
permit us to cover in this course.

\section{And the Rest}

The remaining variations between LGPL and GPL cover the following
conditions:

\begin{itemize}

\item Allowing a licensing ``upgrade'' from LGPL to GPL\@ (in LGPL \S 3)

\item Binary distribution of the library only, covered in LGPL \S 4,
  which is effectively equivalent to LGPL \S 3

\item Creating aggregates of libraries that are not derivative works of
  each other, and distributing them as a unit (in LGPL \S 7)

\end{itemize}


Due to time constraints, we cannot cover these additional terms in detail,
but they are mostly straightforward. The key to understanding LGPL is
understanding the difference between a ``work based on the library'' and a
``work that uses the library.''  Once that distinction is clear, the
remainder of LGPL is close enough to GPL that the concepts discussed in
our more extensive GPL unit can be directly applied.

%%%%%%%%%%%%%%%%%%%%%%%%%%%%%%%%%%%%%%%%%%%%%%%%%%%%%%%%%%%%%%%%%%%%%%%%%%%%%%%
\chapter{Integrating the GPL into Business Practices}

Since GPL'd software is now extremely prevalent through the industry, it
is useful to have some basic knowledge about using GPL'd software in
business and how to build business models around GPL'd software.

\section{Using GPL'd Software In-House}

As discussed in Sections~\ref{GPLs0} and~\ref{GPLs5} of this tutorial,
the GPL only governs the activities of copying, modifying and
distributing software programs that are not governed by the license.
Thus, in FSF's view, simply installing the software on a machine and
using it is not controlled or limited in any way by GPL\@. Using Free
Software in general requires substantially fewer agreements and less
license compliance activity than any known proprietary software.

Even if a company engages heavily in copying the software throughout the
enterprise, such copying is not only permitted by \S\S 1 and 3, but it is
encouraged!  If the company simply deploys unmodified (or even modified)
Free Software throughout the organization for its employees to use, the
obligations under the license are very minimal. Using Free Software has a
substantially lower cost of ownership --- both in licensing fees and in
licensing checking and handling -- than the proprietary software
equivalents.

\section{Business Models}
\label{Business Models}

Using Free Software in house is certainly helpful, but a thriving
market for Free Software-oriented business models also exists. There is the
traditional model of selling copies of Free Software distributions.
Many companies, including IBM and Red Hat, make substantial revenue
from this model. IBM primarily chooses this model because they have
found that for higher-end hardware, the cost of the profit made from
proprietary software licensing fees is negligible. The real profit is
in the hardware, but it is essential that software be stable, reliable
and dependable, and the users be allowed to have unfettered access to
it. Free Software, and GPL'd software in particular (because IBM can
be assured that proprietary versions of the same software will not
exists to compete on their hardware) is the right choice.

Red Hat has actually found that a ``convenience fee'' for Free Software,
when set at a reasonable price (around \$60 or so), can produce some
profit. Even though Red Hat's system is fully downloadable on their
Web site, people still go to local computer stores and buy copies of their
box set, which is simply a printed version of the manual (available under
a Free license as well) and the Free Software system it documents.

\medskip

However, custom support, service, and software improvement contracts
are the most widely used models for GPL'd software. The GPL is
central to their success, because it ensures that the code base
remains common, and that large and small companies are on equal
footing for access to the technology. Consider, for example, the GNU
Compiler Collection (GCC). Cygnus Solutions, a company started in the
early 1990s, was able to grow steadily simply by providing services
for GCC --- mostly consisting of new ports of GCC to different or new,
embedded targets. Eventually, Cygnus was so successful that
it was purchased by Red Hat where it remains a profitable division.

However, there are very small companies like CodeSourcery, as well as
other medium-sized companies like MontaVista and OpenTV that compete in
this space. Because the code-base is protect by GPL, it creates and
demands industry trust. Companies can cooperate on the software and
improve it for everyone. Meanwhile, companies who rely on GCC for their
work are happy to pay for improvements, and for ports to new target
platforms. Nearly all the changes fold back into the standard
versions, and those forks that exist remain freely available.

\medskip

\label{Proprietary Relicensing}

A final common business model that is perhaps the most controversial is
proprietary relicensing of a GPL'd code base. This is only an option for
software in which a particular entity is the sole copyright holder. As
discussed earlier in this tutorial, a copyright holder is permitted under
copyright law to license a software system under her copyright as many
different ways as she likes to as many different parties as she wishes.

Some companies, such as MySQL AB and TrollTech, use this to their
financial advantage with regard to a GPL'd code base. The standard
version is available from the company under the terms of the GPL\@.
However, parties can purchase separate proprietary software licensing for
a fee.

This business model is problematic because it means that the GPL'd code
base must be developed in a somewhat monolithic way, because volunteer
Free Software developers may be reluctant to assign their copyrights to
the company because it will not promise to always and forever license the
software as Free Software. Indeed, the company will surely use such code
contributions in proprietary versions licensed for fees.

\section{Ongoing Compliance}

GPL compliance is in fact a very simple matter -- much simpler than
typical proprietary software agreements and EULAs. Usually, the most
difficult hurdle is changing from a proprietary software mindset to one
that seeks to foster a community of sharing and mutual support. Certainly
complying with the GPL from a users' perspective gives substantially fewer
headaches than proprietary license compliance.

For those who go into the business of distributing {\em modified\\}
versions of GPL'd software, the burden is a bit higher, but not by
much. The glib answer is that by releasing the whole product as Free
Software, it is always easy to comply with the GPL. However,
admittedly to the dismay of FSF, many modern and complex software
systems are built using both proprietary and GPL'd components that are
not legally derivative works of each other. Sometimes, it is easier simply to
improve existing GPL'd application than to start from scratch. In
exchange for that benefit, the license requires that the modifier give
back to the commons that made the work easier in the first place. It is a
reasonable trade-off and a way to help build a better world while also
making a profit.

Note that FSF does provide services to assist companies who need
assistance in complying with the GPL. You can contact FSF's GPL
Compliance Labs at $<$compliance@fsf.org$>$.

If you are particularly interested in matters of GPL compliance, we
recommend the second course in this series, {\em GPL Compliance Case
  Studies and Legal Ethics in Free Software Licensing\/}, in which we
discuss some real GPL violation cases that FSF has worked to resolve.
Consideration of such cases can help give insight on how to handle GPL
compliance in new situations.

\backmatter

\appendix

\chapter{The GNU General Public License}

\begin{center}
{\parindent 0in

Version 2, June 1991

Copyright \copyright\ 1989, 1991 Free Software Foundation, Inc.

\bigskip

59 Temple Place - Suite 330, Boston, MA  02111-1307, USA

\bigskip

Everyone is permitted to copy and distribute verbatim copies
of this license document, but changing it is not allowed.
}
\end{center}

\begin{center}
{\bf\large Preamble}
\end{center}


The licenses for most software are designed to take away your freedom
to share and change it. By contrast, the GNU General Public License is
intended to guarantee your freedom to share and change Free
Software---to make sure the software is free for all its users. This
General Public License applies to most of the Free Software
Foundation's software and to any other program whose authors commit to
using it. (Some other Free Software Foundation software is covered by
the GNU Library General Public License instead.) You can apply it to
your programs, too.

When we speak of Free Software, we are referring to freedom, not price.
Our General Public Licenses are designed to make sure that you have the
freedom to distribute copies of Free Software (and charge for this service
if you wish), that you receive source code or can get it if you want it,
that you can change the software or use pieces of it in new Free programs;
and that you know you can do these things.

To protect your rights, we need to make restrictions that forbid anyone to
deny you these rights or to ask you to surrender the rights. These
restrictions translate to certain responsibilities for you if you
distribute copies of the software, or if you modify it.

For example, if you distribute copies of such a program, whether gratis or
for a fee, you must give the recipients all the rights that you have. You
must make sure that they, too, receive or can get the source code. And
you must show them these terms so they know their rights.

We protect your rights with two steps: (1) copyright the software, and (2)
offer you this license which gives you legal permission to copy,
distribute and/or modify the software.

Also, for each author's protection and ours, we want to make certain that
everyone understands that there is no warranty for this Free Software. If
the software is modified by someone else and passed on, we want its
recipients to know that what they have is not the original, so that any
problems introduced by others will not reflect on the original authors'
reputations.

Finally, any Free program is threatened constantly by software patents.
We wish to avoid the danger that redistributors of a Free program will
individually obtain patent licenses, in effect making the program
proprietary. To prevent this, we have made it clear that any patent must
be licensed for everyone's free use or not licensed at all.

The precise terms and conditions for copying, distribution and
modification follow.

\begin{center}
{\Large \sc Terms and Conditions For Copying, Distribution and
  Modification}
\end{center}


\begin{enumerate}

\addtocounter{enumi}{-1}
\item

This License applies to any program or other work which contains a notice
placed by the copyright holder saying it may be distributed under the
terms of this General Public License. The ``Program,'' below, refers to
any such program or work, and a ``work based on the Program'' means either
the Program or any derivative work under copyright law: that is to say, a
work containing the Program or a portion of it, either verbatim or with
modifications and/or translated into another language. (Hereinafter,
translation is included without limitation in the term ``modification.'')
Each licensee is addressed as ``you.''

Activities other than copying, distribution and modification are not
covered by this License; they are outside its scope. The act of
running the Program is not restricted, and the output from the Program
is covered only if its contents constitute a work based on the
Program (independent of having been made by running the Program).
Whether that is true depends on what the Program does.

\item You may copy and distribute verbatim copies of the Program's source
  code as you receive it, in any medium, provided that you conspicuously
  and appropriately publish on each copy an appropriate copyright notice
  and disclaimer of warranty; keep intact all the notices that refer to
  this License and to the absence of any warranty; and give any other
  recipients of the Program a copy of this License along with the Program.

You may charge a fee for the physical act of transferring a copy, and you
may at your option offer warranty protection in exchange for a fee.

\item

You may modify your copy or copies of the Program or any portion
of it, thus forming a work based on the Program, and copy and
distribute such modifications or work under the terms of Section 1
above, provided that you also meet all of these conditions:

\begin{enumerate}

\item

You must cause the modified files to carry prominent notices stating that
you changed the files and the date of any change.

\item

You must cause any work that you distribute or publish, that in
whole or in part contains or is derived from the Program or any
part thereof, to be licensed as a whole at no charge to all third
parties under the terms of this License.

\item
If the modified program normally reads commands interactively
when run, you must cause it, when started running for such
interactive use in the most ordinary way, to print or display an
announcement including an appropriate copyright notice and a
notice that there is no warranty (or else, saying that you provide
a warranty) and that users may redistribute the program under
these conditions, and telling the user how to view a copy of this
License. (Exception: if the Program itself is interactive but
does not normally print such an announcement, your work based on
the Program is not required to print an announcement.)

\end{enumerate}


These requirements apply to the modified work as a whole. If
identifiable sections of that work are not derived from the Program,
and can be reasonably considered independent and separate works in
themselves, then this License, and its terms, do not apply to those
sections when you distribute them as separate works. But when you
distribute the same sections as part of a whole which is a work based
on the Program, the distribution of the whole must be on the terms of
this License, whose permissions for other licensees extend to the
entire whole, and thus to each and every part regardless of who wrote it.

Thus, it is not the intent of this section to claim rights or contest
your rights to work written entirely by you; rather, the intent is to
exercise the right to control the distribution of derivative or
collective works based on the Program.

In addition, mere aggregation of another work not based on the Program
with the Program (or with a work based on the Program) on a volume of
a storage or distribution medium does not bring the other work under
the scope of this License.

\item
You may copy and distribute the Program (or a work based on it,
under Section 2) in object code or executable form under the terms of
Sections 1 and 2 above provided that you also do one of the following:

\begin{enumerate}

\item

Accompany it with the complete corresponding machine-readable
source code, which must be distributed under the terms of Sections
1 and 2 above on a medium customarily used for software interchange; or,

\item

Accompany it with a written offer, valid for at least three
years, to give any third party, for a charge no more than your
cost of physically performing source distribution, a complete
machine-readable copy of the corresponding source code, to be
distributed under the terms of Sections 1 and 2 above on a medium
customarily used for software interchange; or,

\item

Accompany it with the information you received as to the offer
to distribute corresponding source code. (This alternative is
allowed only for noncommercial distribution and only if you
received the program in object code or executable form with such
an offer, in accord with Subsection b above.)

\end{enumerate}


The source code for a work means the preferred form of the work for
making modifications to it. For an executable work, complete source
code means all the source code for all modules it contains, plus any
associated interface definition files, plus the scripts used to
control compilation and installation of the executable. However, as a
special exception, the source code distributed need not include
anything that is normally distributed (in either source or binary
form) with the major components (compiler, kernel, and so on) of the
operating system on which the executable runs, unless that component
itself accompanies the executable.

If distribution of executable or object code is made by offering
access to copy from a designated place, then offering equivalent
access to copy the source code from the same place counts as
distribution of the source code, even though third parties are not
compelled to copy the source along with the object code.

\item
You may not copy, modify, sublicense, or distribute the Program
except as expressly provided under this License. Any attempt
otherwise to copy, modify, sublicense or distribute the Program is
void, and will automatically terminate your rights under this License.
However, parties who have received copies, or rights, from you under
this License will not have their licenses terminated so long as such
parties remain in full compliance.

\item
You are not required to accept this License, since you have not
signed it. However, nothing else grants you permission to modify or
distribute the Program or its derivative works. These actions are
prohibited by law if you do not accept this License. Therefore, by
modifying or distributing the Program (or any work based on the
Program), you indicate your acceptance of this License to do so, and
all its terms and conditions for copying, distributing or modifying
the Program or works based on it.

\item
Each time you redistribute the Program (or any work based on the
Program), the recipient automatically receives a license from the
original licensor to copy, distribute or modify the Program subject to
these terms and conditions. You may not impose any further
restrictions on the recipients' exercise of the rights granted herein.
You are not responsible for enforcing compliance by third parties to
this License.

\item
If, as a consequence of a court judgment or allegation of patent
infringement or for any other reason (not limited to patent issues),
conditions are imposed on you (whether by court order, agreement or
otherwise) that contradict the conditions of this License, they do not
excuse you from the conditions of this License. If you cannot
distribute so as to satisfy simultaneously your obligations under this
License and any other pertinent obligations, then as a consequence you
may not distribute the Program at all. For example, if a patent
license would not permit royalty-free redistribution of the Program by
all those who receive copies directly or indirectly through you, then
the only way you could satisfy both it and this License would be to
refrain entirely from distribution of the Program.

If any portion of this section is held invalid or unenforceable under
any particular circumstance, the balance of the section is intended to
apply and the section as a whole is intended to apply in other
circumstances.

It is not the purpose of this section to induce you to infringe any
patents or other property right claims or to contest validity of any
such claims; this section has the sole purpose of protecting the
integrity of the Free Software distribution system, which is
implemented by public license practices. Many people have made
generous contributions to the wide range of software distributed
through that system in reliance on consistent application of that
system; it is up to the author/donor to decide if he or she is willing
to distribute software through any other system and a licensee cannot
impose that choice.

This section is intended to make thoroughly clear what is believed to
be a consequence of the rest of this License.

\item
If the distribution and/or use of the Program is restricted in
certain countries either by patents or by copyrighted interfaces, the
original copyright holder who places the Program under this License
may add an explicit geographical distribution limitation excluding
those countries, so that distribution is permitted only in or among
countries not thus excluded. In such case, this License incorporates
the limitation as if written in the body of this License.

\item
The Free Software Foundation may publish revised and/or new versions
of the General Public License from time to time. Such new versions will
be similar in spirit to the present version, but may differ in detail to
address new problems or concerns.

Each version is given a distinguishing version number. If the Program
specifies a version number of this License which applies to it and ``any
later version,'' you have the option of following the terms and conditions
either of that version or of any later version published by the Free
Software Foundation. If the Program does not specify a version number of
this License, you may choose any version ever published by the Free Software
Foundation.

\item
If you wish to incorporate parts of the Program into other free
programs whose distribution conditions are different, write to the author
to ask for permission. For software which is copyrighted by the Free
Software Foundation, write to the Free Software Foundation; we sometimes
make exceptions for this. Our decision will be guided by the two goals
of preserving the free status of all derivatives of our Free Software and
of promoting the sharing and reuse of software generally.

\begin{center}
{\Large\sc
No Warranty
}
\end{center}

\item
{\sc Because the program is licensed free of charge, there is no warranty
for the program, to the extent permitted by applicable law. Except when
otherwise stated in writing the copyright holders and/or other parties
provide the program ``as is'' without warranty of any kind, either expressed
or implied, including, but not limited to, the implied warranties of
merchantability and fitness for a particular purpose. The entire risk as
to the quality and performance of the program is with you. Should the
program prove defective, you assume the cost of all necessary servicing,
repair or correction.}

\item {\sc In no event unless required by applicable law or agreed to
    in writing will any copyright holder, or any other party who may
    modify and/or redistribute the program as permitted above, be
    liable to you for damages, including any general, special,
    incidental or consequential damages arising out of the use or
    inability to use the program (including but not limited to loss of
    data or data being rendered inaccurate or losses sustained by you
    or third parties or a failure of the program to operate with any
    other programs), even if such holder or other party has been
    advised of the possibility of such damages.}

\end{enumerate}


\begin{center}
{\Large\sc End of Terms and Conditions}
\end{center}
\vfill

\pagebreak[4]

\section*{Appendix: How to Apply These Terms to Your New Programs}

If you develop a new program, and you want it to be of the greatest
possible use to the public, the best way to achieve this is to make it
Free Software which everyone can redistribute and change under these
terms.

  To do so, attach the following notices to the program. It is safest to
  attach them to the start of each source file to most effectively convey
  the exclusion of warranty; and each file should have at least the
  ``copyright'' line and a pointer to where the full notice is found.

\begin{quote}
one line to give the program's name and a brief idea of what it does. \\
Copyright (C) yyyy  name of author \\

This program is Free Software; you can redistribute it and/or modify
it under the terms of the GNU General Public License as published by
the Free Software Foundation; either version 2 of the License, or
(at your option) any later version.

This program is distributed in the hope that it will be useful,
but WITHOUT ANY WARRANTY; without even the implied warranty of
MERCHANTABILITY or FITNESS FOR A PARTICULAR PURPOSE. See the
GNU General Public License for more details.

You should have received a copy of the GNU General Public License
along with this program; if not, write to the Free Software
Foundation, Inc., 59 Temple Place - Suite 330, Boston, MA  02111-1307, USA.
\end{quote}

Also add information on how to contact you by electronic and paper mail.

If the program is interactive, make it output a short notice like this
when it starts in an interactive mode:

\begin{quote}
Gnomovision version 69, Copyright (C) yyyy  name of author \\
Gnomovision comes with ABSOLUTELY NO WARRANTY; for details type `show w'. \\
This is Free Software, and you are welcome to redistribute it
under certain conditions; type `show c' for details.
\end{quote}


The hypothetical commands {\tt show w} and {\tt show c} should show the
appropriate parts of the General Public License. Of course, the commands
you use may be called something other than {\tt show w} and {\tt show c};
they could even be mouse-clicks or menu items---whatever suits your
program.

You should also get your employer (if you work as a programmer) or your
school, if any, to sign a ``copyright disclaimer'' for the program, if
necessary. Here is a sample; alter the names:

\begin{quote}
Yoyodyne, Inc., hereby disclaims all copyright interest in the program \\
`Gnomovision' (which makes passes at compilers) written by James Hacker. \\

signature of Ty Coon, 1 April 1989 \\
Ty Coon, President of Vice
\end{quote}


This General Public License does not permit incorporating your program
into proprietary programs. If your program is a subroutine library, you
may consider it more useful to permit linking proprietary applications
with the library. If this is what you want to do, use the GNU Library
General Public License instead of this License.


\chapter{The GNU Lesser General Public License}

\begin{center}
{\parindent 0in

Version 2.1, February 1999

Copyright \copyright\ 1991, 1999 Free Software Foundation, Inc.

\bigskip

59 Temple Place - Suite 330, Boston, MA  02111-1307, USA

\bigskip

Everyone is permitted to copy and distribute verbatim copies
of this license document, but changing it is not allowed.

\bigskip

[This is the first released version of the Lesser GPL. It also counts
 as the successor of the GNU Library Public License version 2, hence
 the version number 2.1.]
}

\end{center}

\begin{center}
{\bf\large Preamble}
\end{center}

The licenses for most software are designed to take away your freedom to
share and change it. By contrast, the GNU General Public Licenses are
intended to guarantee your freedom to share and change Free Software---to
make sure the software is free for all its users.

This license, the Lesser General Public License, applies to some specially
designated software packages---typically libraries---of the Free Software
Foundation and other authors who decide to use it. You can use it too,
but we suggest you first think carefully about whether this license or the
ordinary General Public License is the better strategy to use in any
particular case, based on the explanations below.

When we speak of Free Software, we are referring to freedom of use, not
price. Our General Public Licenses are designed to make sure that you
have the freedom to distribute copies of Free Software (and charge for
this service if you wish); that you receive source code or can get it if
you want it; that you can change the software and use pieces of it in new
Free programs; and that you are informed that you can do these things.

To protect your rights, we need to make restrictions that forbid
distributors to deny you these rights or to ask you to surrender these
rights. These restrictions translate to certain responsibilities for you
if you distribute copies of the library or if you modify it.

For example, if you distribute copies of the library, whether gratis or
for a fee, you must give the recipients all the rights that we gave you.
You must make sure that they, too, receive or can get the source code. If
you link other code with the library, you must provide complete object
files to the recipients, so that they can relink them with the library
after making changes to the library and recompiling it. And you must show
them these terms so they know their rights.

We protect your rights with a two-step method: (1) we copyright the
library, and (2) we offer you this license, which gives you legal
permission to copy, distribute and/or modify the library.

To protect each distributor, we want to make it very clear that there is
no warranty for the Free library. Also, if the library is modified by
someone else and passed on, the recipients should know that what they have
is not the original version, so that the original author's reputation will
not be affected by problems that might be introduced by others.

Finally, software patents pose a constant threat to the existence of any
Free program. We wish to make sure that a company cannot effectively
restrict the users of a Free program by obtaining a restrictive license
from a patent holder. Therefore, we insist that any patent license
obtained for a version of the library must be consistent with the full
freedom of use specified in this license.

Most GNU software, including some libraries, is covered by the ordinary
GNU General Public License. This license, the GNU Lesser General Public
License, applies to certain designated libraries, and is quite different
from the ordinary General Public License. We use this license for certain
libraries in order to permit linking those libraries into non-Free
programs.

When a program is linked with a library, whether statically or using a
shared library, the combination of the two is legally speaking a combined
work, a derivative of the original library. The ordinary General Public
License therefore permits such linking only if the entire combination fits
its criteria of freedom. The Lesser General Public License permits more
lax criteria for linking other code with the library.

We call this license the ``Lesser'' General Public License because it does
Less to protect the user's freedom than the ordinary General Public
License. It also provides other Free Software developers Less of an
advantage over competing non-Free programs. These disadvantages are the
reason we use the ordinary General Public License for many libraries.
However, the Lesser license provides advantages in certain special
circumstances.

For example, on rare occasions, there may be a special need to encourage
the widest possible use of a certain library, so that it becomes a
de-facto standard. To achieve this, non-Free programs must be allowed to
use the library. A more frequent case is that a Free library does the
same job as widely used non-Free libraries. In this case, there is little
to gain by limiting the Free library to Free Software only, so we use the
Lesser General Public License.

In other cases, permission to use a particular library in non-Free
programs enables a greater number of people to use a large body of Free
software. For example, permission to use the GNU C Library in non-Free
programs enables many more people to use the whole GNU operating system,
as well as its variant, the GNU/Linux operating system.

Although the Lesser General Public License is Less protective of the
users' freedom, it does ensure that the user of a program that is linked
with the library has the freedom and the wherewithal to run that program
using a modified version of the library.

The precise terms and conditions for copying, distribution and
modification follow. Pay close attention to the difference between a
``work based on the library'' and a ``work that uses the library.''  The
former contains code derived from the library, whereas the latter must be
combined with the library in order to run.

\begin{center}
{\Large \sc GNU Lesser General Public License} \\
{\Large \sc Terms and Conditions For Copying, Distribution and
  Modification}
\end{center}

\begin{enumerate}

\addtocounter{enumi}{-1}

\item

This License Agreement applies to any software library or other program
which contains a notice placed by the copyright holder or other authorized
party saying it may be distributed under the terms of this Lesser General
Public License (also called ``this License''). Each licensee is addressed
as ``you.''

A ``library'' means a collection of software functions and/or data
prepared so as to be conveniently linked with application programs (which
use some of those functions and data) to form executables.

The ``library,'' below, refers to any such software library or work which
has been distributed under these terms. A ``work based on the library''
means either the library or any derivative work under copyright law: that
is to say, a work containing the library or a portion of it, either
verbatim or with modifications and/or translated straightforwardly into
another language. (Hereinafter, translation is included without
limitation in the term ``modification.'')

``Source code'' for a work means the preferred form of the work for making
modifications to it. For a library, complete source code means all the
source code for all modules it contains, plus any associated interface
definition files, plus the scripts used to control compilation and
installation of the library.

Activities other than copying, distribution and modification are not
covered by this License; they are outside its scope. The act of running a
program using the library is not restricted, and output from such a
program is covered only if its contents constitute a work based on the
library (independent of the use of the library in a tool for writing it).
Whether that is true depends on what the library does and what the program
that uses the library does.
  
\item 

You may copy and distribute verbatim copies of the library's complete
source code as you receive it, in any medium, provided that you
conspicuously and appropriately publish on each copy an appropriate
copyright notice and disclaimer of warranty; keep intact all the notices
that refer to this License and to the absence of any warranty; and
distribute a copy of this License along with the library.

You may charge a fee for the physical act of transferring a copy,
and you may at your option offer warranty protection in exchange for a
fee.

\item

You may modify your copy or copies of the library or any portion of it,
thus forming a work based on the library, and copy and distribute such
modifications or work under the terms of Section 1 above, provided that
you also meet all of these conditions:

\begin{enumerate}

  \item

    The modified work must itself be a software library.

  \item

    You must cause the files modified to carry prominent notices stating
    that you changed the files and the date of any change.

  \item

    You must cause the whole of the work to be licensed at no charge to
    all third parties under the terms of this License.

  \item 
    If a facility in the modified library refers to a function or a table
    of data to be supplied by an application program that uses the
    facility, other than as an argument passed when the facility is
    invoked, then you must make a good faith effort to ensure that, in the
    event an application does not supply such function or table, the
    facility still operates, and performs whatever part of its purpose
    remains meaningful.

(For example, a function in a library to compute square roots has a
purpose that is entirely well-defined independent of the application.
Therefore, Subsection 2d requires that any application-supplied function
or table used by this function must be optional: if the application does
not supply it, the square root function must still compute square roots.)
\end{enumerate}

These requirements apply to the modified work as a whole. If identifiable
sections of that work are not derived from the library, and can be
reasonably considered independent and separate works in themselves, then
this License, and its terms, do not apply to those sections when you
distribute them as separate works. But when you distribute the same
sections as part of a whole which is a work based on the library, the
distribution of the whole must be on the terms of this License, whose
permissions for other licensees extend to the entire whole, and thus to
each and every part regardless of who wrote it.

Thus, it is not the intent of this section to claim rights or contest your
rights to work written entirely by you; rather, the intent is to exercise
the right to control the distribution of derivative or collective works
based on the library.

In addition, mere aggregation of another work not based on the library
with the library (or with a work based on the library) on a volume of a
storage or distribution medium does not bring the other work under the
scope of this License.

\item

You may opt to apply the terms of the ordinary GNU General Public License
instead of this License to a given copy of the library. To do this, you
must alter all the notices that refer to this License, so that they refer
to the ordinary GNU General Public License version 2, instead of to this
License. (If a newer version than version 2 of the ordinary GNU General
Public License has appeared, then you can specify that version instead if
you wish.)  Do not make any other change in these notices.

Once this change is made in a given copy, it is irreversible for that
copy, so the ordinary GNU General Public License applies to all subsequent
copies and derivative works made from that copy.

This option is useful when you wish to copy part of the code of the
library into a program that is not a library.

\item

You may copy and distribute the library (or a portion or derivative of it,
under Section 2) in object code or executable form under the terms of
Sections 1 and 2 above provided that you accompany it with the complete
corresponding machine-readable source code, which must be distributed
under the terms of Sections 1 and 2 above on a medium customarily used for
software interchange.

If distribution of object code is made by offering access to copy from a
designated place, then offering equivalent access to copy the source code
from the same place satisfies the requirement to distribute the source
code, even though third parties are not compelled to copy the source along
with the object code.

\item

A program that contains no derivative of any portion of the library, but
is designed to work with the library by being compiled or linked with it,
is called a ``work that uses the library.''  Such a work, in isolation, is
not a derivative work of the library, and therefore falls outside the
scope of this License.

However, linking a ``work that uses the library'' with the library creates
an executable that is a derivative of the library (because it contains
portions of the library), rather than a ``work that uses the library.''
The executable is therefore covered by this License. Section 6 states
terms for distribution of such executables.

When a ``work that uses the library'' uses material from a header file
that is part of the library, the object code for the work may be a
derivative work of the library even though the source code is not.
Whether this is true is especially significant if the work can be linked
without the library, or if the work is itself a library. The threshold
for this to be true is not precisely defined by law.

If such an object file uses only numerical parameters, data structure
layouts and accessors, and small macros and small inline functions (ten
lines or less in length), then the use of the object file is unrestricted,
regardless of whether it is legally a derivative work. (Executables
containing this object code plus portions of the library will still fall
under Section 6.)

Otherwise, if the work is a derivative of the library, you may distribute
the object code for the work under the terms of Section 6. Any
executables containing that work also fall under Section 6, whether or not
they are linked directly with the library itself.

\item

As an exception to the Sections above, you may also combine or link a
``work that uses the library'' with the library to produce a work
containing portions of the library, and distribute that work under terms
of your choice, provided that the terms permit modification of the work
for the customer's own use and reverse engineering for debugging such
modifications.

You must give prominent notice with each copy of the work that the library
is used in it and that the library and its use are covered by this
License. You must supply a copy of this License. If the work during
execution displays copyright notices, you must include the copyright
notice for the library among them, as well as a reference directing the
user to the copy of this License. Also, you must do one of these things:

\begin{enumerate}

  \item

    Accompany the work with the complete corresponding machine-readable
    source code for the library including whatever changes were used in
    the work (which must be distributed under Sections 1 and 2 above);
    and, if the work is an executable linked with the library, with the
    complete machine-readable ``work that uses the library,'' as object
    code and/or source code, so that the user can modify the library and
    then relink to produce a modified executable containing the modified
    library. (It is understood that the user who changes the contents of
    definitions files in the library will not necessarily be able to
    recompile the application to use the modified definitions.)

  \item

    Use a suitable shared library mechanism for linking with the library.
    A suitable mechanism is one that (1) uses at run time a copy of the
    library already present on the user's computer system, rather than
    copying library functions into the executable, and (2) will operate
    properly with a modified version of the library, if the user installs
    one, as long as the modified version is interface-compatible with the
    version that the work was made with.

  \item

    Accompany the work with a written offer, valid for at least three
    years, to give the same user the materials specified in Subsection 6a,
    above, for a charge no more than the cost of performing this
    distribution.

  \item

    If distribution of the work is made by offering access to copy from a
    designated place, offer equivalent access to copy the above specified
    materials from the same place.

  \item

    Verify that the user has already received a copy of these materials or
    that you have already sent this user a copy.
\end{enumerate}

For an executable, the required form of the ``work that uses the library''
must include any data and utility programs needed for reproducing the
executable from it. However, as a special exception, the materials to be
distributed need not include anything that is normally distributed (in
either source or binary form) with the major components (compiler, kernel,
and so on) of the operating system on which the executable runs, unless
that component itself accompanies the executable.

It may happen that this requirement contradicts the license restrictions
of other proprietary libraries that do not normally accompany the
operating system. Such a contradiction means you cannot use both them and
the library together in an executable that you distribute.

\item

You may place library facilities that are a work based on the library
side-by-side in a single library together with other library facilities
not covered by this License, and distribute such a combined library,
provided that the separate distribution of the work based on the library
and of the other library facilities is otherwise permitted, and provided
that you do these two things:

\begin{enumerate}

   \item

     Accompany the combined library with a copy of the same work based on
     the library, uncombined with any other library facilities. This must
     be distributed under the terms of the Sections above.

   \item

     Give prominent notice with the combined library of the fact that part
     of it is a work based on the library, and explaining where to find
     the accompanying uncombined form of the same work.
\end{enumerate}

\item

  You may not copy, modify, sublicense, link with, or distribute the
  library except as expressly provided under this License. Any attempt
  otherwise to copy, modify, sublicense, link with, or distribute the
  library is void, and will automatically terminate your rights under this
  License. However, parties who have received copies, or rights, from you
  under this License will not have their licenses terminated so long as
  such parties remain in full compliance.

\item  

  You are not required to accept this License, since you have not signed
  it. However, nothing else grants you permission to modify or distribute
  the library or its derivative works. These actions are prohibited by
  law if you do not accept this License. Therefore, by modifying or
  distributing the library (or any work based on the library), you
  indicate your acceptance of this License to do so, and all its terms and
  conditions for copying, distributing or modifying the library or works
  based on it.

\item

  Each time you redistribute the library (or any work based on the
  library), the recipient automatically receives a license from the
  original licensor to copy, distribute, link with or modify the library
  subject to these terms and conditions. You may not impose any further
  restrictions on the recipients' exercise of the rights granted herein.
  You are not responsible for enforcing compliance by third parties with
  this License.

\item

  If, as a consequence of a court judgment or allegation of patent
  infringement or for any other reason (not limited to patent issues),
  conditions are imposed on you (whether by court order, agreement or
  otherwise) that contradict the conditions of this License, they do not
  excuse you from the conditions of this License. If you cannot
  distribute so as to satisfy simultaneously your obligations under this
  License and any other pertinent obligations, then as a consequence you
  may not distribute the library at all. For example, if a patent license
  would not permit royalty-free redistribution of the library by all those
  who receive copies directly or indirectly through you, then the only way
  you could satisfy both it and this License would be to refrain entirely
  from distribution of the library.

  If any portion of this section is held invalid or unenforceable under
  any particular circumstance, the balance of the section is intended to
  apply, and the section as a whole is intended to apply in other
  circumstances.

  It is not the purpose of this section to induce you to infringe any
  patents or other property right claims or to contest validity of any
  such claims; this section has the sole purpose of protecting the
  integrity of the Free Software distribution system which is implemented
  by public license practices. Many people have made generous
  contributions to the wide range of software distributed through that
  system in reliance on consistent application of that system; it is up to
  the author/donor to decide if he or she is willing to distribute
  software through any other system and a licensee cannot impose that
  choice.

  This section is intended to make thoroughly clear what is believed to be
  a consequence of the rest of this License.


% \pagebreak[4]


\item 

  If the distribution and/or use of the library is restricted in certain
  countries either by patents or by copyrighted interfaces, the original
  copyright holder who places the library under this License may add an
  explicit geographical distribution limitation excluding those countries,
  so that distribution is permitted only in or among countries not thus
  excluded. In such case, this License incorporates the limitation as if
  written in the body of this License.

\item 

  The Free Software Foundation may publish revised and/or new versions of
  the Lesser General Public License from time to time. Such new versions
  will be similar in spirit to the present version, but may differ in
  detail to address new problems or concerns.

  Each version is given a distinguishing version number. If the library
  specifies a version number of this License which applies to it and ``any
  later version,'' you have the option of following the terms and
  conditions either of that version or of any later version published by
  the Free Software Foundation. If the library does not specify a license
  version number, you may choose any version ever published by the Free
  Software Foundation.


\item
  

  If you wish to incorporate parts of the library into other Free programs
  whose distribution conditions are incompatible with these, write to the
  author to ask for permission. For software which is copyrighted by the
  Free Software Foundation, write to the Free Software Foundation; we
  sometimes make exceptions for this. Our decision will be guided by the
  two goals of preserving the Free status of all derivatives of our Free
  software and of promoting the sharing and reuse of software generally.


\begin{center}
{\Large\sc
No Warranty
}
\end{center}

\item

{\sc Because the library is licensed free of charge, there is no
warranty for the library, to the extent permitted by applicable law.
Except when otherwise stated in writing the copyright holders and/or
other parties provide the library ``as is'' without warranty of any
kind, either expressed or implied, including, but not limited to, the
implied warranties of merchantability and fitness for a particular
purpose. The entire risk as to the quality and performance of the
library is with you. should the library prove defective, you assume
the cost of all necessary servicing, repair or correction.}

% \pagebreak[4]

\item

{\sc In no event unless required by applicable law or agreed to in writing
  will any copyright holder, or any other party who may modify and/or
  redistribute the library as permitted above, be liable to you for
  damages, including any general, special, incidental or consequential
  damages arising out of the use or inability to use the library
  (including but not limited to loss of data or data being rendered
  inaccurate or losses sustained by you or third parties or a failure of
  the library to operate with any other software), even if such holder or
  other party has been advised of the possibility of such damages.}

\end{enumerate}

\begin{center}
{\Large\sc End of Terms and Conditions}
\end{center}
\vfill

\pagebreak[4]

\section*{How to Apply These Terms to Your New Libraries}
           
If you develop a new library, and you want it to be of the greatest
possible use to the public, we recommend making it Free Software that
everyone can redistribute and change. You can do so by permitting
redistribution under these terms (or, alternatively, under the terms of
the ordinary General Public License).

To apply these terms, attach the following notices to the library. It is
safest to attach them to the start of each source file to most effectively
convey the exclusion of warranty; and each file should have at least the
``copyright'' line and a pointer to where the full notice is found.

\begin{quote}
one line to give the library's name and a brief idea of what it does. \\
Copyright (C) year  name of author \\

This library is Free Software; you can redistribute it and/or modify it
under the terms of the GNU Lesser General Public License as published by
the Free Software Foundation; either version 2.1 of the License, or (at
your option) any later version.

This library is distributed in the hope that it will be useful, but
WITHOUT ANY WARRANTY; without even the implied warranty of MERCHANTABILITY
or FITNESS FOR A PARTICULAR PURPOSE. See the GNU Lesser General Public
License for more details.

You should have received a copy of the GNU Lesser General Public License
along with this library; if not, write to the Free Software Foundation,
Inc., 59 Temple Place, Suite 330, Boston, MA 02111-1307 USA
\end{quote}

Also add information on how to contact you by electronic and paper mail.

You should also get your employer (if you work as a programmer) or your
school, if any, to sign a ``copyright disclaimer'' for the library, if
necessary. Here is a sample; alter the names:

\begin{quote}
Yoyodyne, Inc., hereby disclaims all copyright interest in the program \\
`Gnomovision' (which makes passes at compilers) written by James Hacker. \\

signature of Ty Coon, 1 April 1990 \\
Ty Coon, President of Vice
\end{quote}


\chapter{The Affero General Public License}

\begin{center}
{\parindent 0in

Version 1, March 2002

Copyright \copyright\ 2002 Affero, Inc.

\bigskip

510 Third Street - Suite 225, San Francisco, CA 94107, USA

\bigskip

This license is a modified version of the GNU General Public License
copyright (C) 1989, 1991 Free Software Foundation, Inc. made with
their permission. Section 2(d) has been added to cover use of software
over a computer network.

Everyone is permitted to copy and distribute verbatim copies
of this license document, but changing it is not allowed.
}
\end{center}

\begin{center}
{\bf\large Preamble}
\end{center}



The licenses for most software are designed to take away your freedom
to share and change it. By contrast, the Affero General Public License
is intended to guarantee your freedom to share and change free
software--to make sure the software is free for all its users. This
Public License applies to most of Affero's software and to any other
program whose authors commit to using it. (Some other Affero software
is covered by the GNU Library General Public License instead.) You can
apply it to your programs, too.


When we speak of free software, we are referring to freedom, not price. This General Public License is designed to make sure that you have the freedom to distribute copies of free software (and charge for this service if you wish), that you receive source code or can get it if you want it, that you can change the software or use pieces of it in new free programs; and that you know you can do these things.

To protect your rights, we need to make restrictions that forbid anyone to deny you these rights or to ask you to surrender the rights. These restrictions translate to certain responsibilities for you if you distribute copies of the software, or if you modify it.

For example, if you distribute copies of such a program, whether gratis or for a fee, you must give the recipients all the rights that you have. You must make sure that they, too, receive or can get the source code. And you must show them these terms so they know their rights.

We protect your rights with two steps: (1) copyright the software, and (2) offer you this license which gives you legal permission to copy, distribute and/or modify the software.

Also, for each author's protection and ours, we want to make certain that everyone understands that there is no warranty for this free software. If the software is modified by someone else and passed on, we want its recipients to know that what they have is not the original, so that any problems introduced by others will not reflect on the original authors' reputations.

Finally, any free program is threatened constantly by software patents. We wish to avoid the danger that redistributors of a free program will individually obtain patent licenses, in effect making the program proprietary. To prevent this, we have made it clear that any patent must be licensed for everyone's free use or not licensed at all.

The precise terms and conditions for copying, distribution and modification follow.

\begin{center}
{\Large \sc Terms and Conditions For Copying, Distribution and
  Modification}
\end{center}


\begin{enumerate}

\addtocounter{enumi}{-1}
\item

This License applies to any program or other work which contains a
notice placed by the copyright holder saying it may be distributed
under the terms of this Affero General Public License.  The
``Program'', below, refers to any such program or work, and a ``work
based on the Program'' means either the Program or any derivative work
under copyright law: that is to say, a work containing the Program or
a portion of it, either verbatim or with modifications and/or
translated into another language.  (Hereinafter, translation is
included without limitation in the term ``modification''.)  Each
licensee is addressed as ``you''.

Activities other than copying, distribution and modification are not
covered by this License; they are outside its scope.  The act of
running the Program is not restricted, and the output from the Program
is covered only if its contents constitute a work based on the
Program (independent of having been made by running the Program).
Whether that is true depends on what the Program does.

\item You may copy and distribute verbatim copies of the Program's source
  code as you receive it, in any medium, provided that you conspicuously
  and appropriately publish on each copy an appropriate copyright notice
  and disclaimer of warranty; keep intact all the notices that refer to
  this License and to the absence of any warranty; and give any other
  recipients of the Program a copy of this License along with the Program.

You may charge a fee for the physical act of transferring a copy, and you
may at your option offer warranty protection in exchange for a fee.

\item

You may modify your copy or copies of the Program or any portion
of it, thus forming a work based on the Program, and copy and
distribute such modifications or work under the terms of Section 1
above, provided that you also meet all of these conditions:

\begin{enumerate}

\item

You must cause the modified files to carry prominent notices stating that
you changed the files and the date of any change.

\item

You must cause any work that you distribute or publish, that in
whole or in part contains or is derived from the Program or any
part thereof, to be licensed as a whole at no charge to all third
parties under the terms of this License.

\item
If the modified program normally reads commands interactively
when run, you must cause it, when started running for such
interactive use in the most ordinary way, to print or display an
announcement including an appropriate copyright notice and a
notice that there is no warranty (or else, saying that you provide
a warranty) and that users may redistribute the program under
these conditions, and telling the user how to view a copy of this
License.  (Exception: if the Program itself is interactive but
does not normally print such an announcement, your work based on
the Program is not required to print an announcement.)

\item
\textbf{If the Program as you received it is intended to interact with users
through a computer network and if, in the version you received, any
user interacting with the Program was given the opportunity to request
transmission to that user of the Program's complete source code, you
must not remove that facility from your modified version of the
Program or work based on the Program, and must offer an equivalent
opportunity for all users interacting with your Program through a
computer network to request immediate transmission by HTTP of the
complete source code of your modified version or other derivative
work.}

\end{enumerate}


These requirements apply to the modified work as a whole.  If
identifiable sections of that work are not derived from the Program,
and can be reasonably considered independent and separate works in
themselves, then this License, and its terms, do not apply to those
sections when you distribute them as separate works.  But when you
distribute the same sections as part of a whole which is a work based
on the Program, the distribution of the whole must be on the terms of
this License, whose permissions for other licensees extend to the
entire whole, and thus to each and every part regardless of who wrote it.

Thus, it is not the intent of this section to claim rights or contest
your rights to work written entirely by you; rather, the intent is to
exercise the right to control the distribution of derivative or
collective works based on the Program.

In addition, mere aggregation of another work not based on the Program
with the Program (or with a work based on the Program) on a volume of
a storage or distribution medium does not bring the other work under
the scope of this License.

\item
You may copy and distribute the Program (or a work based on it,
under Section 2) in object code or executable form under the terms of
Sections 1 and 2 above provided that you also do one of the following:

\begin{enumerate}

\item

Accompany it with the complete corresponding machine-readable
source code, which must be distributed under the terms of Sections
1 and 2 above on a medium customarily used for software interchange; or,

\item

Accompany it with a written offer, valid for at least three
years, to give any third party, for a charge no more than your
cost of physically performing source distribution, a complete
machine-readable copy of the corresponding source code, to be
distributed under the terms of Sections 1 and 2 above on a medium
customarily used for software interchange; or,

\item

Accompany it with the information you received as to the offer
to distribute corresponding source code.  (This alternative is
allowed only for noncommercial distribution and only if you
received the program in object code or executable form with such
an offer, in accord with Subsection b above.)

\end{enumerate}


The source code for a work means the preferred form of the work for
making modifications to it.  For an executable work, complete source
code means all the source code for all modules it contains, plus any
associated interface definition files, plus the scripts used to
control compilation and installation of the executable.  However, as a
special exception, the source code distributed need not include
anything that is normally distributed (in either source or binary
form) with the major components (compiler, kernel, and so on) of the
operating system on which the executable runs, unless that component
itself accompanies the executable.

If distribution of executable or object code is made by offering
access to copy from a designated place, then offering equivalent
access to copy the source code from the same place counts as
distribution of the source code, even though third parties are not
compelled to copy the source along with the object code.

\item
You may not copy, modify, sublicense, or distribute the Program
except as expressly provided under this License.  Any attempt
otherwise to copy, modify, sublicense or distribute the Program is
void, and will automatically terminate your rights under this License.
However, parties who have received copies, or rights, from you under
this License will not have their licenses terminated so long as such
parties remain in full compliance.

\item
You are not required to accept this License, since you have not
signed it.  However, nothing else grants you permission to modify or
distribute the Program or its derivative works.  These actions are
prohibited by law if you do not accept this License.  Therefore, by
modifying or distributing the Program (or any work based on the
Program), you indicate your acceptance of this License to do so, and
all its terms and conditions for copying, distributing or modifying
the Program or works based on it.

\item
Each time you redistribute the Program (or any work based on the
Program), the recipient automatically receives a license from the
original licensor to copy, distribute or modify the Program subject to
these terms and conditions.  You may not impose any further
restrictions on the recipients' exercise of the rights granted herein.
You are not responsible for enforcing compliance by third parties to
this License.

\item
If, as a consequence of a court judgment or allegation of patent
infringement or for any other reason (not limited to patent issues),
conditions are imposed on you (whether by court order, agreement or
otherwise) that contradict the conditions of this License, they do not
excuse you from the conditions of this License.  If you cannot
distribute so as to satisfy simultaneously your obligations under this
License and any other pertinent obligations, then as a consequence you
may not distribute the Program at all.  For example, if a patent
license would not permit royalty-free redistribution of the Program by
all those who receive copies directly or indirectly through you, then
the only way you could satisfy both it and this License would be to
refrain entirely from distribution of the Program.

If any portion of this section is held invalid or unenforceable under
any particular circumstance, the balance of the section is intended to
apply and the section as a whole is intended to apply in other
circumstances.

It is not the purpose of this section to induce you to infringe any
patents or other property right claims or to contest validity of any
such claims; this section has the sole purpose of protecting the
integrity of the free software distribution system, which is
implemented by public license practices.  Many people have made
generous contributions to the wide range of software distributed
through that system in reliance on consistent application of that
system; it is up to the author/donor to decide if he or she is willing
to distribute software through any other system and a licensee cannot
impose that choice.

This section is intended to make thoroughly clear what is believed to
be a consequence of the rest of this License.

\item
If the distribution and/or use of the Program is restricted in
certain countries either by patents or by copyrighted interfaces, the
original copyright holder who places the Program under this License
may add an explicit geographical distribution limitation excluding
those countries, so that distribution is permitted only in or among
countries not thus excluded.  In such case, this License incorporates
the limitation as if written in the body of this License.

\item
\textbf{Affero Inc. may publish revised and/or new versions of the Affero
General Public License from time to time. Such new versions will be
similar in spirit to the present version, but may differ in detail to
address new problems or concerns.}

\textbf{Each version is given a distinguishing version number. If the Program
specifies a version number of this License which applies to it and
``any later version'', you have the option of following the terms and
conditions either of that version or of any later version published by
Affero, Inc. If the Program does not specify a version number of this
License, you may choose any version ever published by Affero, Inc.}

\textbf{You may also choose to redistribute modified versions of this program
under any version of the Free Software Foundation's GNU General Public
License version 3 or higher, so long as that version of the GNU GPL
includes terms and conditions substantially equivalent to those of
this license.}

\item
If you wish to incorporate parts of the Program into other free
programs whose distribution conditions are different, write to the
author to ask for permission. For software which is copyrighted by
Affero, Inc., write to us; we sometimes make exceptions for this. Our
decision will be guided by the two goals of preserving the free status
of all derivatives of our free software and of promoting the sharing
and reuse of software generally.

\begin{center}
{\Large\sc
No Warranty
}
\end{center}

\item
{\sc Because the program is licensed free of charge, there is no warranty
for the program, to the extent permitted by applicable law.  Except when
otherwise stated in writing the copyright holders and/or other parties
provide the program ``as is'' without warranty of any kind, either expressed
or implied, including, but not limited to, the implied warranties of
merchantability and fitness for a particular purpose.  The entire risk as
to the quality and performance of the program is with you.  Should the
program prove defective, you assume the cost of all necessary servicing,
repair or correction.}

\item
{\sc In no event unless required by applicable law or agreed to in writing
will any copyright holder, or any other party who may modify and/or
redistribute the program as permitted above, be liable to you for damages,
including any general, special, incidental or consequential damages arising
out of the use or inability to use the program (including but not limited
to loss of data or data being rendered inaccurate or losses sustained by
you or third parties or a failure of the program to operate with any other
programs), even if such holder or other party has been advised of the
possibility of such damages.}

\end{enumerate}

That's all there is to it!

\chapter{GPL Version 3: Background to Adoption}

\textbf{\textit{\large{by Richard Stallman and Eben Moglen}}}

\smallskip

The GNU General Public License (``the GPL'') has remained
unmodified, at version level 2, since 1991.  This is extraordinary
longevity for any widely-employed legal instrument.  The durability of
the GPL is even more surprising when one takes into account the
differences between the free software movement at the time of version
2's release and the situation prevailing in 2005.

Richard M. Stallman, founder of the free software movement and author
of the GNU GPL, released version 2 in 1991 after taking legal advice
and collecting developer opinion concerning version 1 of the license,
which had been in use since 1985.  There was no formal public comment
process and no significant interim transition period.  The Free
Software Foundation immediately relicensed the components of the GNU
Project, which comprised the largest then-existing collection of
copyleft software assets.  In Finland, Linus Torvalds adopted GPL
Version 2 for his operating system kernel, called Linux.

That was then, and this is now.  The GPL is employed by tens of
thousands of software projects around the world, of which the Free
Software Foundation's GNU system is a tiny fraction.  The GNU system,
when combined with Linus Torvalds' Linux---which has evolved into a
flexible, highly-portable, industry-leading operating system kernel
--- along with Samba, MySQL, and other GPL'd programs, offers superior
reliability and adaptability to Microsoft's operating systems, at
nominal cost.  GPL'd software runs on or is embedded in devices
ranging from cellphones, PDAs and home networking appliances to
mainframes and supercomputing clusters.  Independent software
developers around the world, as well as every large corporate IT buyer
and seller, and a surprisingly large proportion of individual users,
interact with the GPL.

During the period since 1991, of course, there has developed a
profusion of free software licenses.  But not in the area covered by
the GPL.  The ``share and share alike'' or ``copyleft'' aspect of the
GPL is its most important functional characteristic, and those who
want to use a copyleft license for software overwhelmingly use the GPL
rather than inventing their own.

Updating the GPL is therefore a very different task in 2005 than it
was in 1991.  The substantive reasons for revision, and the likely
nature of those changes, are subject matter for another essay.  At
present we would like to concentrate on the institutional, procedural
aspects of changing the license.  Those are complicated by the fact
that the GPL serves four distinct purposes.

\section*{The GPL is a Worldwide Copyright License}

As a legal document, the GPL serves a purpose that most legal drafters
would do anything possible to avoid: it licenses copyrighted material
for modification and redistribution in every one of the world's
systems of copyright law.  In general, publishers don't use worldwide
copyright licenses; for each system in which their works are
distributed, licensing arrangements tailored to local legal
requirements are used.  Publishers rarely license redistribution of
modified or derivative works; when they do so, those licenses are
tailored to the specific setting, factual and legal.  But free
software requires legal arrangements that permit copyrighted works to
follow arbitrary trajectories, in both geographic and genetic terms.
Modified versions of free software works are distributed from hand to
hand across borders in a pattern that no copyright holder could
possibly trace.

GPL version 2 performed the task of globalization relatively well,
because its design was elegantly limited to a minimum set of copyright
principles that signatories to the Berne Convention must offer, in one
form or another, in their national legislation.  But GPL2 was a
license constructed by one US layman and his lawyers, largely
concerned with US law.  To the extent possible, and without any
fundamental changes, GPL3 should ease internationalization
difficulties, more fully approximating the otherwise unsought ideal of
the global copyright license.

\section*{The GPL is the Code of Conduct for Free Software Distributors}

Beyond the legal permission that the GPL extends to those who wish to
copy, modify, and share free software, the GPL also embodies a code of
industry conduct with respect to the practices by which free software
is distributed.  Section 3, which explains how to make source code
available as required under the license, affects product packaging
decisions for those who embed free software in appliances, as well as
those who distribute software collections that include both free and
unfree software.  Section 7, which concerns the effect of licenses,
judgments, and other compulsory legal interventions incompatible with
the GPL on the behavior of software distributors, affects patent
licensing arrangements in connection with industry standards.  And so
on, through a range of interactions between the requirements of the
license and evolving practices in the vending of both hardware and
software.

The Free Software Foundation, through its maintenance and enforcement
of the GPL, has contributed to the evolution of industry behavior
patterns beyond its influence as a maker of software.  In revising the
GPL, the Foundation is inevitably engaged in altering the rules of the
road for enterprises and market participants of many different kinds,
with different fundamental interests and radically different levels of
market power.  The process of drafting and adopting changes to the
license must thus approximate standard-setting, or ``best practices''
definition, as well as copyright license drafting.

\section*{The GPL is the Constitution of the Free Software Movement}

The Free Software Foundation has never been reluctant to point out
that its goals are primarily social and political, not technical or
economic.  The Foundation believes that free software---that is,
software that can be freely studied, copied, modified, reused,
redistributed and shared by its users---is the only ethically
satisfactory form of software development, as free and open scientific
research is the only ethically satisfactory context for the conduct of
mathematics, physics, or biology.  The Foundation, and those who
support its broader work, regard free software as an essential step in
a social movement for freer access to knowledge, freer access to
facilities of communication, and a more deeply participatory culture,
open to human beings with less regard to existing distributions of
wealth and social power.  The free software movement has taken
advantage of the social conditions of its time to found its program on
the creation of vast new wealth, through new systems of cooperation,
which can in turn be shared in order to further the creation of new
wealth, in a positive feedback loop.

This program is not, of course, universally shared by all the parties
who benefit from the exploitation of the new wealth created by free
software.  The free software movement has never objected to the
indirect benefits accruing to those who differ from the movement's
goals: one of the powerful lessons the movement has learned from
previous aspects of the long-duration Western movement for freedom of
expression is the value of working with, rather than against,
conventional economic interests and concerns.  But the movement's own
goals cannot be subordinated to the economic interests of our friends
and allies in industry, let alone those who occasionally contribute
solely for reasons of their own.  Changes to the GPL, for whatever
reason they are undertaken, must not undermine the underlying movement
for freer exchange of knowledge.  To the extent that the movement has
identified technological or legal measures likely to be harmful to
freedom, such as ``trusted computing'' or a broadening of the scope of
patent law, the GPL needs to address those issues from a perspective
of political principle and the needs of the movement, not from primary
regard for the industrial or commercial consequences.

\section*{The GPL is the Literary Work of Richard M.\ Stallman}

Some copyright licenses are no doubt known, in the restricted circle
of one firm or law office, as the achievement of a single author's
acumen or insight.  But it is safe to say that there is no other
copyright license in the world that is so strongly identified with the
achievements, and the philosophy, of a single public figure.  Mr.\
Stallman remains the GPL's author, with as much right to preserve its
integrity as a work representative of his intentions as any other
author or creator.  Under his guidance, the Free Software Foundation,
which holds the copyright of the GPL, will coordinate and direct the
process of its modification.

\section*{Conclusion}

The GPL serves, and must continue to serve, multiple purposes.  Those
purposes are fundamentally diverse, and they inevitably conflict.
Development of GPL version 3 has been an ongoing process within the
Free Software Foundation; we, along with our colleagues, have never
stopped considering possible modifications.  We have consulted,
formally and informally, a very broad array of participants in the
free software community, from industry, the academy, and the garage.
Those conversations have occurred in many countries and several
languages, over almost two decades, as the technology of software
development and distribution changed around us.

When a GPLv3 discussion draft is released, the pace of that
conversation will change, as a particular proposal becomes the
centerpiece.  The Foundation will, before it emits a first discussion
draft, publicize the process by which it intends to gather opinion and
suggestions.  The Free Software Foundation recognizes that the
reversioning of the GPL is a crucial moment in the evolution of the
free software community, and the Foundation intends to meet its
responsibilities to the makers, distributors and users of free
software.  In doing so, we hope to hear all relevant points of view,
and to make decisions that reflect the many disparate purposes that
the license must serve.  Our primary concern remains, as it has been
from the beginning, the creation and protection of freedom.  We
recognize that the best protection of freedom is a growing and vital
community of the free.  We will use the process of public discussion
of GPL3 drafts to support and nurture the community of the free.
Proprietary culture imposes both technology and license terms; free
software means allowing people to understand, experiment and modify
software, as well as getting involved in the discussion of license
terms, so that everyone's ideas can contribute to the common good, and
the development of each contributes to the development of all.

% =====================================================================
% END OF FIRST DAY SEMINAR SECTION
% =====================================================================

\end{document}
